\begin{proof}[Proof of \cref{obj:prop:capacity-identification}]
Throughout, conditional probabilities on a zero-mass conditioning event are interpreted by the totalized convention as zero; on positive-mass conditioning events they are the usual ratios.  Write
\[
\mathsf{Cell}_x=\{X=x\},\qquad
\mathsf{Arm}_{x,z}=\{X=x,Z=z\}.
\]
Cell probabilities are nonnegative because \(p_x=P(\mathsf{Cell}_x)\).  Now fix \(x\) with \(p_x>0\).  For the conditional instrument propensity \(\pi(x)\), the totalized convention is the usual ratio on this cell, so
\[
\pi(x)=\frac{P(\mathsf{Arm}_{x,1})}{P(\mathsf{Cell}_x)},
\qquad
P(\mathsf{Arm}_{x,1})=\pi(x)p_x .
\]
By \cref{obj:ass:instrument-overlap},
\[
0<\varepsilon_Z\le \pi(x)\le 1-\varepsilon_Z<1.
\]
Hence \(P(\mathsf{Arm}_{x,1})>0\).  Since the two instrument fibers partition the cell,
\[
p_x=P(\mathsf{Cell}_x)=P(\mathsf{Arm}_{x,0})+P(\mathsf{Arm}_{x,1}),
\]
and therefore
\[
P(\mathsf{Arm}_{x,0})=p_x-P(\mathsf{Arm}_{x,1})=(1-\pi(x))p_x>0.
\]
Thus \(P(\mathsf{Arm}_{x,z})>0\) for both \(z\in\{0,1\}\).
% lean: arm_cell_positive

For supported \(x\), the observed lower-arm contrast pulls back along the observed-data law to
\[
\ell_i(x)
=
P(D=0,S=1,Y=i\mid X=x,Z=0)
-
P(D=0,S=1,Y=i\mid X=x,Z=1).
\]
Similarly,
\[
h_j(x)
=
P(D=1,S=1,Y=j\mid X=x,Z=1)
-
P(D=1,S=1,Y=j\mid X=x,Z=0).
\]
The observed-data law is the pushforward of the latent law by the observed datum map, whose coordinate events are measurable; hence these are exactly the corresponding conditional probabilities under the latent probability space.
% lean: observable_lower_pullback
% lean: observable_upper_pullback

Using \cref{obj:ass:iv-independence,obj:ass:treatment-consistency,obj:ass:selection-exclusion-consistency,obj:ass:outcome-exclusion-consistency}, for each \(z\in\{0,1\}\),
\[
P(D=0,S=1,Y=i\mid X=x,Z=z)
=
P(D_z=0,S_0=1,Y_0=i\mid X=x),
\]
and
\[
P(D=1,S=1,Y=j\mid X=x,Z=z)
=
P(D_z=1,S_1=1,Y_1=j\mid X=x).
\]
Substitution gives
\[
\ell_i(x)
=
P(D_0=0,S_0=1,Y_0=i\mid X=x)
-
P(D_1=0,S_0=1,Y_0=i\mid X=x),
\]
and
\[
h_j(x)
=
P(D_1=1,S_1=1,Y_1=j\mid X=x)
-
P(D_0=1,S_1=1,Y_1=j\mid X=x).
\]
% lean: selectedOutcome_arm_identified

By \cref{obj:ass:no-defiers}, \(D_1\ge D_0\), so
\[
\{D_0=0,S_0=1,Y_0=i\}
=
\{D_1=0,S_0=1,Y_0=i\}\,\dot\cup\,
\{C,S_0=1,Y_0=i\}
\]
up to null sets.  Therefore
\[
\ell_i(x)=P(Y_0=i,S_0=1,C\mid X=x).
\]
% lean: lower_noDefier_contrast

The same monotonicity gives
\[
\{D_1=1,S_1=1,Y_1=j\}
=
\{D_0=1,S_1=1,Y_1=j\}\,\dot\cup\,
\{C,S_1=1,Y_1=j\},
\]
and hence
\[
h_j(x)=P(Y_1=j,S_1=1,C\mid X=x).
\]
% lean: upper_noDefier_contrast

These identities also show nonnegativity in supported cells.  If \(p_x=0\), then \(P(\mathsf{Arm}_{x,z})=0\) for both \(z\), so every conditional term in the observable contrasts is zero and the corresponding capacities are nonnegative.  Thus the observed-data law is a probability law and its observable capacity contrasts are valid, so it is compatible with the maintained capacity domain.
% lean: capacity_identification

Because \(\mathcal Y=\{0,\ldots,K-1\}\) is finite, summing the lower selected-complier identities over outcome levels gives
\[
q_0(x)=\sum_{i\in\mathcal Y}\ell_i(x)
=
P(S_0=1,C\mid X=x).
\]
% lean: sum_lowerLatentMass

Applying the same finite partition to the upper selected-complier identities gives
\[
q_1(x)=\sum_{j\in\mathcal Y}h_j(x)
=
P(S_1=1,C\mid X=x).
\]
% lean: sum_upperLatentMass

By \cref{obj:def:observable-capacities},
\[
\Delta q(x)=q_1(x)-q_0(x)
=
P(S_1=1,C\mid X=x)-P(S_0=1,C\mid X=x).
\]
% lean: capacity_identification

It remains to identify \(m(x)\).  Define the finite restricted measure \(\nu_x\) by
\[
\nu_x(A)=P(A\cap C\cap\mathsf{Cell}_x)
\]
for every measurable event \(A\).  If \(\nu_x(\Omega)=0\), all selected-complier masses in the cell are zero.  Otherwise, \cref{obj:ass:weak-selection-monotonicity} applies on this restricted cell.  If \(d(x)=+1\), then \(S_0\le S_1\) \(\nu_x\)-almost surely, so
\[
P(S_0=1,S_1=1,C\mid X=x)=P(S_0=1,C\mid X=x),
\]
and
\[
P(S_0=1,C\mid X=x)\le P(S_1=1,C\mid X=x).
\]
If \(d(x)=-1\), then \(S_1\le S_0\) \(\nu_x\)-almost surely, so
\[
P(S_0=1,S_1=1,C\mid X=x)=P(S_1=1,C\mid X=x),
\]
and
\[
P(S_1=1,C\mid X=x)\le P(S_0=1,C\mid X=x).
\]
Combining these alternatives with \(m(x)=\min\{q_0(x),q_1(x)\}\) from \cref{obj:def:observable-capacities} yields
\[
m(x)=P(S_0=1,S_1=1,C\mid X=x).
\]
% lean: selection_monotone_cell

The same two alternatives determine every strict sign of the gap.  If \(d(x)=-1\), then \(q_1(x)\le q_0(x)\), so \(\Delta q(x)\le0\).  Hence \(\Delta q(x)>0\) forces \(d(x)=+1\).  If \(d(x)=+1\), then \(q_0(x)\le q_1(x)\), so \(\Delta q(x)\ge0\).  Hence \(\Delta q(x)<0\) forces \(d(x)=-1\).
% lean: capacity_identification

Finally suppose \(p_x>0\) and \(\Delta q(x)=0\).  The gap identity gives
\[
P(S_1=1,C\mid X=x)=P(S_0=1,C\mid X=x).
\]
Equivalently, \(\nu_x(\{S_1=1\})=\nu_x(\{S_0=1\})\).  If \(\nu_x(\Omega)=0\), both direction inequalities are vacuous in the complier cell.  If \(\nu_x(\Omega)>0\), \cref{obj:ass:weak-selection-monotonicity} gives one almost-sure inclusion between \(\{S_0=1\}\) and \(\{S_1=1\}\) under \(\nu_x\); equality of their finite measures upgrades that inclusion to
\[
\nu_x(\{S_0\ne S_1\})=0.
\]
Thus, in either case, both inequalities \(S_0\le S_1\) and \(S_1\le S_0\) hold \(\nu_x\)-almost surely in the cell.  Define the updated direction maps by
\[
d^{+x}(u)=
\begin{cases}
+1,&u=x,\\
d(u),&u\ne x,
\end{cases}
\qquad
 d^{-x}(u)=
\begin{cases}
-1,&u=x,\\
d(u),&u\ne x.
\end{cases}
\]
Away from \(x\), the original weak-monotonicity condition is unchanged; at \(x\), the preceding alternatives supply either direction.  Therefore the same latent law, with direction map \(d^{+x}\), witnesses observational admissibility of \(+1\), and the same latent law, with direction map \(d^{-x}\), witnesses observational admissibility of \(-1\).  The observed-data law is unchanged in both witnesses.
% lean: weakMonotone_update_of_equal
\end{proof}

\begin{proof}[Proof of \cref{obj:prop:tie-face-collapse}]
Fix a cell \(x\) with \(p_x>0\), and let \(c\) be the observable capacity vector from \(P_{\mathrm{obs}}\). By \cref{obj:prop:capacity-identification}, \(c\) is valid, so
\[
\ell_i(x)\ge 0,\qquad h_j(x)\ge 0
\]
for all \(i,j\in\mathcal Y\), and
\[
\Delta q(x)
=
P(S_1=1,C\mid X=x)-P(S_0=1,C\mid X=x).
\]
Also, by \cref{obj:def:observable-capacities},
\[
q_0(x)=\sum_i \ell_i(x),\qquad
q_1(x)=\sum_j h_j(x),\qquad
\Delta q(x)=q_1(x)-q_0(x),\qquad
m(x)=\min\{q_0(x),q_1(x)\}.
\]
For a coupling array \(\gamma_x\), write
\[
\operatorname{row}_i(\gamma_x)=\sum_j \gamma_{ij,x},
\qquad
\operatorname{col}_j(\gamma_x)=\sum_i \gamma_{ij,x},
\qquad
\operatorname{tot}(\gamma_x)=\sum_i\sum_j\gamma_{ij,x}.
\]
These totals satisfy
\[
\sum_i\operatorname{row}_i(\gamma_x)=\operatorname{tot}(\gamma_x),
\qquad
\sum_j\operatorname{col}_j(\gamma_x)=\operatorname{tot}(\gamma_x).
\]
% lean: capacity_identification

\begin{enumerate}
\item We first prove the face equivalence. If \(\gamma_x\in\Gamma_x^\ast\), then \cref{obj:def:partial-transport-polytope} gives
\[
\operatorname{row}_i(\gamma_x)\le \ell_i(x),\qquad
\operatorname{col}_j(\gamma_x)\le h_j(x),\qquad
\operatorname{tot}(\gamma_x)=m(x).
\]
When \(q_0(x)\le q_1(x)\),
\[
\sum_i\{\ell_i(x)-\operatorname{row}_i(\gamma_x)\}
=
q_0(x)-m(x)
=
0.
\]
Each summand is nonnegative, hence
\[
\operatorname{row}_i(\gamma_x)=\ell_i(x)\qquad\text{for every }i.
\]
Similarly, when \(q_1(x)\le q_0(x)\),
\[
\operatorname{col}_j(\gamma_x)=h_j(x)\qquad\text{for every }j.
\]

By the cellwise output of \cref{obj:def:threshold-flow-construction}, choose a feasible exact-mass coupling
\[
\gamma_x^{\mathrm L}\in \Gamma_x^\ast .
\]
Assume
\[
\Gamma_x^{+}=\Gamma_x^{-},
\qquad
\Gamma_x^{-}=\Gamma_x^{0}.
\]
If \(q_0(x)\le q_1(x)\), the preceding row-exactness argument puts
\(\gamma_x^{\mathrm L}\) in \(\Gamma_x^{+}\). The displayed face equalities then put
\(\gamma_x^{\mathrm L}\) in \(\Gamma_x^0\). Therefore
\[
q_0(x)=\operatorname{tot}(\gamma_x^{\mathrm L})=q_1(x),
\]
using exact rows for the first equality and exact columns for the second. If \(q_1(x)\le q_0(x)\), the same argument with columns first puts
\(\gamma_x^{\mathrm L}\) in \(\Gamma_x^{-}\), then in \(\Gamma_x^0\), and again gives
\[
q_0(x)=\operatorname{tot}(\gamma_x^{\mathrm L})=q_1(x).
\]
Thus \(\Delta q(x)=0\).

Conversely, suppose \(\Delta q(x)=0\). Then
\[
q_0(x)=q_1(x)=m(x).
\]
If \(\gamma_x\in\Gamma_x^{+}\), then its exact row constraints give
\[
\operatorname{tot}(\gamma_x)
=
\sum_i\operatorname{row}_i(\gamma_x)
=
\sum_i\ell_i(x)
=
q_0(x)
=
m(x),
\]
so \(\gamma_x\in\Gamma_x^\ast\). Since \(q_1(x)\le q_0(x)\), the column-exactness argument above gives
\[
\operatorname{col}_j(\gamma_x)=h_j(x)\qquad\text{for every }j,
\]
and hence \(\gamma_x\in\Gamma_x^{-}\). The reverse inclusion follows symmetrically from exact columns and row-exactness. Therefore
\[
\Gamma_x^{+}=\Gamma_x^{-}.
\]
The same reasoning sends every element of \(\Gamma_x^{-}\) to \(\Gamma_x^0\), while \(\Gamma_x^0\subseteq\Gamma_x^{-}\) is immediate from the defining constraints. Hence
\[
\Gamma_x^{-}=\Gamma_x^0.
\]
This proves
\[
\bigl(\Gamma_x^{+}=\Gamma_x^{-}\ \text{and}\ \Gamma_x^{-}=\Gamma_x^{0}\bigr)
\Longleftrightarrow
\Delta q(x)=0.
\]
% lean: polytope_tie_iff_gap_zero

\item Now assume \(\Delta q(x)=0\). Since \(p_x>0\), conditional probabilities given \(X=x\) use the denominator \(p_x\). The displayed identity from \cref{obj:prop:capacity-identification} yields
\[
P(S_1=1,C\mid X=x)=P(S_0=1,C\mid X=x).
\]
If \(P(C\mid X=x)=0\), then both
\[
P(S_0=S_1,C\mid X=x)
\quad\text{and}\quad
P(C\mid X=x)
\]
are zero. If \(P(C\mid X=x)>0\), then \cref{obj:ass:weak-selection-monotonicity} gives one of the two almost-sure orders on the complier cell:
\[
1\{S_0=1\}\le 1\{S_1=1\}
\quad\text{or}\quad
1\{S_1=1\}\le 1\{S_0=1\}.
\]
The two indicators have equal conditional integrals by the preceding display, so the ordered indicators are equal almost surely on \(\{C,X=x\}\). Therefore
\[
P(S_0=S_1,\ C\mid X=x)=P(C\mid X=x).
\]
% lean: equalSelectionComplierMass_of_zero_gap

\item Under the same zero-gap condition, the argument in the first step also gives
\[
\Gamma_x^\ast=\Gamma_x^0.
\]
Indeed, every \(\gamma_x\in\Gamma_x^\ast\) has exact rows and exact columns because
\(q_0(x)=q_1(x)=m(x)\), and every \(\gamma_x\in\Gamma_x^0\) has bounded rows, bounded columns, and total mass \(m(x)\). Combining this equality with the already proved
\[
\Gamma_x^+=\Gamma_x^-,
\qquad
\Gamma_x^-=\Gamma_x^0,
\]
and with \(\Gamma_x=\Gamma_x^\ast\) from \cref{obj:def:partial-transport-polytope}, gives
\[
\Gamma_x=\Gamma_x^{+},
\qquad
\Gamma_x=\Gamma_x^{-},
\qquad
\Gamma_x=\Gamma_x^{0}.
\]
% lean: branchFree_eq_tie_of_gap_zero

\item Finally fix a threshold \(t\in\mathcal Y\). The prefix sums \(\ell_{\le t}(x)\) and \(h_{\le t}(x)\), and the upper tail \(h_{>t}(x)\), are those of \cref{obj:def:threshold-cuts}. The lower tail appearing in the statement is
\[
\ell_{>t}(x):=\sum_{\substack{i\in\mathcal Y\\ t<i}}\ell_i(x).
\]
Because \(\mathcal Y\) is finite and each index lies in exactly one of the two sets \(\{i:i\le t\}\) and \(\{i:t<i\}\),
\[
\ell_{\le t}(x)+\ell_{>t}(x)
=
\sum_{\substack{i\in\mathcal Y\\ i\le t}}\ell_i(x)
+
\sum_{\substack{i\in\mathcal Y\\ t<i}}\ell_i(x)
=
\sum_{i\in\mathcal Y}\ell_i(x)
=
q_0(x).
\]
The same finite partition gives
\[
h_{\le t}(x)+h_{>t}(x)
=
\sum_{\substack{j\in\mathcal Y\\ j\le t}}h_j(x)
+
\sum_{\substack{j\in\mathcal Y\\ t<j}}h_j(x)
=
\sum_{j\in\mathcal Y}h_j(x)
=
q_1(x).
\]
When \(\Delta q(x)=0\), \(q_0(x)=q_1(x)\), and hence
\[
\ell_{\le t}(x)-h_{\le t}(x)
=
\bigl(q_0(x)-\ell_{>t}(x)\bigr)
-
\bigl(q_1(x)-h_{>t}(x)\bigr)
=
h_{>t}(x)-\ell_{>t}(x).
\]
% lean: prefix_difference_eq_tail_difference_of_gap_zero
\end{enumerate}
The four displayed conclusions are exactly the asserted cellwise conclusions for every \(x\) with \(p_x>0\). % lean: tie_face_collapse
\end{proof}

\begin{proof}[Proof of \cref{obj:thm:full-law-endpoint-attainment}]
Let \(c\) be the observable capacity vector from \cref{obj:def:observable-capacities}, and let
\[
p_x=P_{\mathrm{obs}}(X=x).
\]
By \cref{obj:prop:capacity-identification}, the capacity vector \(c\) is valid, the weights satisfy \(p_x\ge 0\), and, for every supported cell \(p_x>0\),
\[
\ell_i(x)=P(Y_0=i,S_0=1,C\mid X=x),
\qquad
h_j(x)=P(Y_1=j,S_1=1,C\mid X=x),
\]
\[
\Delta q(x)=P(S_1=1,C\mid X=x)-P(S_0=1,C\mid X=x),
\qquad
m(x)=P(S_0=S_1=1,C\mid X=x).
\]
It also gives the direction implications
\[
\Delta q(x)>0\Rightarrow d(x)=+1,
\qquad
\Delta q(x)<0\Rightarrow d(x)=-1.
\]
The capacity-identification step applies on supported cells.  On cells with \(p_x=0\), the observed cell event has zero probability.  Hence, for each instrument arm \(z\),
\[
P_{\mathrm{obs}}(X=x,Z=z)=0.
\]
With the paper's zero-denominator convention for conditional probabilities, the two conditional probabilities entering every lower-arm contrast and every upper-arm contrast are therefore zero.  Thus
\[
\ell_i(x)=0,\qquad h_j(x)=0
\]
for every \(i,j\).  Since
\[
q_0(x)=\sum_i\ell_i(x),\qquad
q_1(x)=\sum_j h_j(x),
\qquad
m(x)=\min\{q_0(x),q_1(x)\},
\]
it follows that
\[
q_0(x)=q_1(x)=m(x)=0.
\]
% lean: observableCapacity_eq_zero_of_p_eq_zero

Choose as baseline the original compatible law generating \(P_{\mathrm{obs}}\).  The following cell quantities use the finite-table totalization convention: if \(p_x>0\), they are the ordinary conditional probabilities given \(X=x\); if \(p_x=0\), they are the corresponding zero-cell baseline table entries.  Write the baseline complier mass as
\[
\bar c_C(x)=P(C\mid X=x),
\]
and write its never-taker and always-taker components as
\[
\mathsf N_x(s_0,s_1,y_0,y_1)
=
P(D_0=D_1=0,S_0=s_0,S_1=s_1,Y_0=y_0,Y_1=y_1\mid X=x),
\]
\[
\mathsf A_x(s_0,s_1,y_0,y_1)
=
P(D_0=D_1=1,S_0=s_0,S_1=s_1,Y_0=y_0,Y_1=y_1\mid X=x).
\]
For \(p_x>0\), the selected-complier masses are bounded by the complier mass, hence
\[
q_0(x)\le \bar c_C(x),
\qquad
q_1(x)\le \bar c_C(x),
\qquad
\max\{q_0(x),q_1(x)\}\le \bar c_C(x).
\]
For \(p_x=0\), both capacity totals are zero, so the same inequality holds.  Thus the baseline is compatible with \(P_{\mathrm{obs}}\) and \(c\), has conditional complier mass at least \(\max\{q_0(x),q_1(x)\}\) in every cell, and carries the displayed never-taker and always-taker components. % lean: originalCompatibleBaseline

Since \(K\ge 3\), fix the reference outcome
\[
\bar y=0\in\mathcal Y.
\]
For each cell, apply the lower and upper threshold-flow constructions from \cref{obj:def:threshold-flow-construction}.  Let
\[
\gamma^L_x(i,j)
\quad\text{and}\quad
\gamma^U_x(i,j)
\]
be the two survivor-complier couplings.  They satisfy, for \(\bullet\in\{L,U\}\),
\[
\gamma^\bullet_x(i,j)\ge 0,\qquad
\sum_j\gamma^\bullet_x(i,j)\le \ell_i(x),\qquad
\sum_i\gamma^\bullet_x(i,j)\le h_j(x),\qquad
\sum_{i,j}\gamma^\bullet_x(i,j)=m(x),
\]
and their benefit masses are
\[
\sum_{i<j}\gamma^L_x(i,j)=B_L(x),
\qquad
\sum_{i<j}\gamma^U_x(i,j)=B_U(x).
\]
% lean: thresholdFlow_spec

For each \(\bullet\in\{L,U\}\), define the one-sided selected-complier residuals by
\[
\rho^{0,\bullet}_{x,i}
=
\begin{cases}
\ell_i(x)-\sum_j\gamma^\bullet_x(i,j),& \Delta q(x)<0,\\
0,& \Delta q(x)\ge 0,
\end{cases}
\]
\[
\rho^{1,\bullet}_{x,j}
=
\begin{cases}
h_j(x)-\sum_i\gamma^\bullet_x(i,j),& \Delta q(x)>0,\\
0,& \Delta q(x)\le 0,
\end{cases}
\]
and the never-selected complier mass by
\[
\nu_x=\bar c_C(x)-\max\{q_0(x),q_1(x)\}.
\]
Then, for every cell and every \(\bullet\in\{L,U\}\),
\[
\sum_j\gamma^\bullet_x(i,j)+\rho^{0,\bullet}_{x,i}=\ell_i(x),
\qquad
\sum_i\gamma^\bullet_x(i,j)+\rho^{1,\bullet}_{x,j}=h_j(x).
\]
Indeed, when \(\Delta q(x)<0\) or \(\Delta q(x)>0\), the corresponding identity is the definition of the residual.  When \(\Delta q(x)\ge0\), the row upper bounds have total \(q_0(x)=m(x)\), so every row bound binds; when \(\Delta q(x)\le0\), the column upper bounds have total \(q_1(x)=m(x)\), so every column bound binds. % lean: thresholdLatentCompletion_selected_margins

Build, for each \(\bullet\in\{L,U\}\), a cell table
\[
\mathfrak t^\bullet_x(d_0,d_1,s_0,s_1,y_0,y_1),
\qquad
(d_0,d_1,s_0,s_1)\in\{0,1\}^4,\quad y_0,y_1\in\mathcal Y.
\]
For \(p_x>0\), set
\[
\mathfrak t^\bullet_x(d_0,d_1,s_0,s_1,y_0,y_1)
=
\begin{cases}
\mathsf N_x(s_0,s_1,y_0,y_1),& d_0=d_1=0,\\
\gamma^\bullet_x(y_0,y_1),& d_0=0,\ d_1=1,\ s_0=s_1=1,\\
\rho^{0,\bullet}_{x,y_0},& d_0=0,\ d_1=1,\ s_0=1,\ s_1=0,\ y_1=\bar y,\\
\rho^{1,\bullet}_{x,y_1},& d_0=0,\ d_1=1,\ s_0=0,\ s_1=1,\ y_0=\bar y,\\
\nu_x,& d_0=0,\ d_1=1,\ s_0=s_1=0,\ y_0=y_1=\bar y,\\
\mathsf A_x(s_0,s_1,y_0,y_1),& d_0=d_1=1,\\
0,& \text{otherwise.}
\end{cases}
\]
For \(p_x=0\), set
\[
\mathfrak t^\bullet_x(d_0,d_1,s_0,s_1,y_0,y_1)
=
\begin{cases}
1,& d_0=d_1=s_0=s_1=0,\ y_0=y_1=\bar y,\\
0,& \text{otherwise.}
\end{cases}
\]
The nonnegative entries sum to one in every cell: in supported cells the survivor, one-sided, and never-selected complier entries sum to \(\bar c_C(x)\), while \cref{obj:ass:no-defiers} sets the defier stratum \(\{D_0=1,D_1=0\}\) to zero, so the never-taker, complier, and always-taker components exhaust the principal-stratum partition; in unsupported cells the table is a point mass. % lean: normalizedCompletionCellTable_normalized

Let \(P^\bullet\) be the law obtained by drawing \(X\) with probabilities \(p_x\), drawing \(Z\) conditionally on \(X=x\) with the same instrument propensity \(\pi(x)\) as \(P_{\mathrm{obs}}\), and drawing \((D_0,D_1,S_0,S_1,Y_0,Y_1)\) conditionally on \(X=x\) according to \(\mathfrak t^\bullet_x\), independently of \(Z\) given \(X\).  The associated slate sets
\[
D=D_Z,\qquad S=S_D,\qquad Y=Y_D\ \text{on }\{S=1\}.
\]
The conditional product construction gives conditional instrument independence. % lean: canonicalThresholdCandidate_ivIndependence
The displayed definitions of \(D\), \(S\), and \(Y\) give treatment consistency, selection consistency, and outcome consistency. % lean: canonicalThresholdCandidate_consistency
The entry with \(D_0=1,D_1=0\) has zero mass in every cell of \(\mathfrak t^\bullet_x\); hence the normalized cell table has no defier mass, and the induced slate satisfies treatment monotonicity. % lean: normalizedCompletionCellTable_noDefiers, canonicalThresholdCandidate_noDefiers

The weak selection direction is also preserved.  If \(p_x>0\) and \(d(x)=+1\), the implication above gives \(\Delta q(x)<0\) as impossible, hence
\[
\rho^{0,\bullet}_{x,i}=0\quad\text{for all }i,
\]
so the complier stratum with \((S_0,S_1)=(1,0)\) has zero mass.  If \(p_x>0\) and \(d(x)=-1\), then \(\Delta q(x)>0\) is excluded and
\[
\rho^{1,\bullet}_{x,j}=0\quad\text{for all }j,
\]
so the complier stratum with \((S_0,S_1)=(0,1)\) has zero mass.  Unsupported cells have \(p_x=0\), so they contribute no conditional positive-mass case. % lean: canonicalThresholdCandidate_weakSelectionMonotonicity

The observed law is unchanged.  For supported cells, the selected-complier margins displayed above give exactly the observable capacities \(\ell_i(x)\) and \(h_j(x)\); the never-taker and always-taker margins are the baseline ones; and the complier total is
\[
\sum_{i,j}\gamma^\bullet_x(i,j)
+\sum_i\rho^{0,\bullet}_{x,i}
+\sum_j\rho^{1,\bullet}_{x,j}
+\nu_x
=
\bar c_C(x).
\]
For \(z\in\{0,1\}\), put
\[
\varpi_z(x)=
\begin{cases}
1-\pi(x),& z=0,\\
\pi(x),& z=1.
\end{cases}
\]
\Cref{obj:ass:instrument-overlap} gives \(\varpi_z(x)>0\) whenever \(p_x>0\).  For treatment state \(\mathfrak d\in\{0,1\}\), selection state \(\mathfrak s\in\{0,1\}\), and reported outcome \(\upsilon\in\{\varnothing\}\cup\mathcal Y\), define
\[
R_{\mathfrak d}
=
\begin{cases}
\varnothing,& S_{\mathfrak d}=0,\\
Y_{\mathfrak d},& S_{\mathfrak d}=1,
\end{cases}
\]
and define the original latent tail margin
\[
\lambda_x(z,\mathfrak d,\mathfrak s,\upsilon)
=
P\!\left(D_z=\mathfrak d,\ S_{\mathfrak d}=\mathfrak s,\ R_{\mathfrak d}=\upsilon\mid X=x\right).
\]
For the constructed table, define the decoded observed margin
\[
\Lambda^\bullet_x(z,\mathfrak d,\mathfrak s,\upsilon)
=
\sum_{d_0,d_1,s_0,s_1,y_0,y_1}
\mathbf 1\!
\left\{
\begin{array}{c}
 d_z=\mathfrak d,\ s_{\mathfrak d}=\mathfrak s,\\
 \upsilon=\varnothing\ \text{if }s_{\mathfrak d}=0,\quad
 \upsilon=y_{\mathfrak d}\ \text{if }s_{\mathfrak d}=1
\end{array}
\right\}
\mathfrak t^\bullet_x(d_0,d_1,s_0,s_1,y_0,y_1).
\]
For selected reported outcomes, the table definitions and the displayed selected-complier margins give
\[
\Lambda^\bullet_x(0,0,1,i)
=
\sum_{s_1,y_1}\mathsf N_x(1,s_1,i,y_1)+\ell_i(x),
\]
\[
\Lambda^\bullet_x(0,1,1,j)
=
\sum_{s_0,y_0}\mathsf A_x(s_0,1,y_0,j),
\]
\[
\Lambda^\bullet_x(1,0,1,i)
=
\sum_{s_1,y_1}\mathsf N_x(1,s_1,i,y_1),
\]
and
\[
\Lambda^\bullet_x(1,1,1,j)
=
\sum_{s_0,y_0}\mathsf A_x(s_0,1,y_0,j)+h_j(x).
\]
The corresponding original latent margins have the same four values: \cref{obj:prop:capacity-identification} identifies \(\ell_i(x)\) and \(h_j(x)\) with the selected complier outcome masses on supported cells, while \cref{obj:ass:no-defiers} makes the defier terms in the mixed treatment strata equal to zero.  Hence
\[
\Lambda^\bullet_x(z,\mathfrak d,1,i)
=
\lambda_x(z,\mathfrak d,1,i)
\]
for every selected reported outcome \(i\).  The invalid observed-outcome margins also agree:
\[
\Lambda^\bullet_x(z,\mathfrak d,0,i)=\lambda_x(z,\mathfrak d,0,i)=0,
\qquad
\Lambda^\bullet_x(z,\mathfrak d,1,\varnothing)=\lambda_x(z,\mathfrak d,1,\varnothing)=0.
\]
It remains only the unselected reported outcome.  The constructed table has the same conditional mass for each principal treatment event as the baseline, so
\[
\Lambda^\bullet_x(z,\mathfrak d,0,\varnothing)
+
\sum_i\Lambda^\bullet_x(z,\mathfrak d,1,i)
=
P(D_z=\mathfrak d\mid X=x).
\]
The original latent tails form the same finite partition,
\[
\lambda_x(z,\mathfrak d,0,\varnothing)
+
\sum_i\lambda_x(z,\mathfrak d,1,i)
=
P(D_z=\mathfrak d\mid X=x).
\]
Together with the selected-margin identities, this gives
\[
\Lambda^\bullet_x(z,\mathfrak d,0,\varnothing)
=
\lambda_x(z,\mathfrak d,0,\varnothing).
\]
Thus, for every supported cell and every observable singleton,
\[
P^\bullet(X=x,Z=z,D=\mathfrak d,S=\mathfrak s,Y=\upsilon)
=
p_x\varpi_z(x)\Lambda^\bullet_x(z,\mathfrak d,\mathfrak s,\upsilon)
=
p_x\varpi_z(x)\lambda_x(z,\mathfrak d,\mathfrak s,\upsilon)
=
P(X=x,Z=z,D=\mathfrak d,S=\mathfrak s,Y=\upsilon).
\]
The first equality is the decoded-table construction; the last equality uses \cref{obj:ass:iv-independence,obj:ass:treatment-consistency,obj:ass:selection-exclusion-consistency,obj:ass:outcome-exclusion-consistency} and the positive arm mass supplied by \cref{obj:ass:instrument-overlap}.  Cells with \(p_x=0\) have zero observed mass under both laws.  Hence
\[
P^L_{\mathrm{obs}}=P_{\mathrm{obs}},
\qquad
P^U_{\mathrm{obs}}=P_{\mathrm{obs}}.
\]
% lean: canonicalThresholdCandidate_completion_observedLaw_eq
Because the observed law is the same, the inequalities in \cref{obj:ass:instrument-overlap} and the mass condition in \cref{obj:ass:positive-aggregate-survivors} transfer unchanged. % lean: instrumentOverlap_of_observedLaw_eq, positiveAggregateSurvivors_of_observedLaw_eq

It remains to compute the target value.  In supported cells, the survivor coupling of \(P^\bullet\) is \(\gamma^\bullet_x\); in unsupported cells its contribution is multiplied by \(p_x=0\).  Hence
\[
\sum_x p_x\sum_{i,j}\gamma^\bullet_x(i,j)
=
\sum_x p_xm(x)
=
M.
\]
The benefit numerators are
\[
\sum_x p_x\sum_{i<j}\gamma^L_x(i,j)
=
\sum_x p_xB_L(x),
\qquad
\sum_x p_x\sum_{i<j}\gamma^U_x(i,j)
=
\sum_x p_xB_U(x).
\]
Since \(M>0\), division by \(M\) and the endpoint formula in \cref{obj:def:threshold-flow-construction} give
\[
P^L(Y_1>Y_0\mid S_1=S_0=1,C)=\theta_L,
\qquad
P^U(Y_1>Y_0\mid S_1=S_0=1,C)=\theta_U.
\]
Thus \(P^L\) and \(P^U\) are endpoint witnesses. % lean: canonicalThresholdCandidate_benefitProbabilityOf

Finally, the realization statement follows cell by cell.  If \(p_x>0\), the conditional table of \(P^L\) equals the lower completion components just defined, and the conditional table of \(P^U\) equals the upper completion components. % lean: canonicalThresholdCandidate_realizes_of_pos
If \(p_x=0\), the zero-cell baseline convention gives
\[
\ell_i(x)=0,\qquad h_j(x)=0,\qquad \bar c_C(x)=0,
\]
and all baseline never-taker and always-taker component masses in that cell are zero.  Therefore the lower and upper threshold-flow completions are both the zero completion,
\[
\gamma^L_x(i,j)=\gamma^U_x(i,j)=0,\qquad
\rho^{0,L}_{x,i}=\rho^{0,U}_{x,i}=0,\qquad
\rho^{1,L}_{x,j}=\rho^{1,U}_{x,j}=0,
\qquad
\nu_x=0,
\]
for all \(i,j\).  The realization identities in a zero-probability cell are evaluated under the same zero conditioning-cell mass convention: every conditional baseline component used by the completion is zero, and the zero-flow rule sets the lower and upper threshold completions to the zero array.  Thus the survivor, one-sided selected-complier, never-selected, never-taker, and always-taker component functionals of \(P^L\) and \(P^U\) are the zero functionals displayed above. % lean: originalCompatibleBaseline_data_zero, thresholdFlow_isZero_of_data_zero, canonicalThresholdCandidate_realizes_of_zero
Hence \(P^L\) realizes the lower threshold latent completion and \(P^U\) realizes the upper threshold latent completion for every covariate cell. % lean: canonicalThresholdCandidate_realizes_of_pos, canonicalThresholdCandidate_realizes_of_zero
\end{proof}

\begin{proof}[Proof of \cref{obj:thm:no-selection-reduction}]
\begin{enumerate}
\item
Fix a covariate cell \(x\) with \(p_x>0\), and form the observable capacities \(c\) from \(P_{\mathrm{obs}}\).  By \cref{obj:prop:capacity-identification}, these capacities are valid and, in this supported cell,
\[
\Delta q(x)
=
P(S_1=1,C\mid X=x)-P(S_0=1,C\mid X=x).
\]
The no-selection submodel implies the two selected-complier masses equal the complier mass:
\[
P(S_0=1,C\mid X=x)=P(C\mid X=x),
\qquad
P(S_1=1,C\mid X=x)=P(C\mid X=x).
\]
Indeed, if \(P(C\mid X=x)>0\), then
\[
P(S_0=S_1=1\mid C,X=x)=1,
\]
and the event \(\{S_0=S_1=1,C\}\) is contained in each of
\(\{S_0=1,C\}\) and \(\{S_1=1,C\}\), while both are contained in \(C\).
Thus both conditional probabilities are squeezed to \(P(C\mid X=x)\).
If \(P(C\mid X=x)=0\), the same containment in \(C\) gives both selected-complier masses equal to zero.  Substituting these two identities in the displayed formula for \(\Delta q(x)\) gives
\[
\Delta q(x)=0 .
\]
% lean: selectedComplierMass_eq_complierMass_of_noSelection, capacity_identification

\item
It remains to identify the cellwise feasible set.  Since \(\Delta q(x)=q_1(x)-q_0(x)\), the zero-gap identity gives
\[
q_0(x)=q_1(x).
\]
The exact survivor-complier mass is
\[
m(x)=\min\{q_0(x),q_1(x)\}=q_0(x)=q_1(x).
\]
For any \(\gamma_x\in\Gamma_x\), \cref{obj:def:partial-transport-polytope} gives
\[
\sum_j\gamma_{ij,x}\le \ell_i(x),
\qquad
\sum_i\gamma_{ij,x}\le h_j(x),
\qquad
\sum_{i,j}\gamma_{ij,x}=m(x).
\]
Summing the row inequalities gives
\[
m(x)=\sum_{i,j}\gamma_{ij,x}
=\sum_i\sum_j\gamma_{ij,x}
\le \sum_i \ell_i(x)=q_0(x)=m(x),
\]
so every row inequality binds.  The same argument with columns gives
\[
m(x)=\sum_{i,j}\gamma_{ij,x}
=\sum_j\sum_i\gamma_{ij,x}
\le \sum_j h_j(x)=q_1(x)=m(x),
\]
so every column inequality binds.  Hence \(\gamma_x\in\Gamma_x^{0}\).  Conversely, every \(\gamma_x\in\Gamma_x^{0}\) has both exact margins, and therefore has total mass \(q_0(x)=q_1(x)=m(x)\) and satisfies the defining inequalities of \(\Gamma_x\).  Thus
\[
\Gamma_x=\Gamma_x^{0}.
\]
% lean: branchFree_eq_tie_of_gap_zero

\item
Now assume \(m(x)>0\).  By \cref{obj:def:threshold-cuts},
\[
B_L(x)
=
\max\left\{
0,\
\max_{t\in\mathcal Y}\bigl[\ell_{\le t}(x)-h_{\le t}(x)\bigr]
+\min\{\Delta q(x),0\}
\right\}.
\]
Using \(\Delta q(x)=0\), this becomes
\[
B_L(x)
=
\max\left\{
0,\
\max_{t\in\mathcal Y}\bigl[\ell_{\le t}(x)-h_{\le t}(x)\bigr]
\right\}.
\]
Division by the positive number \(m(x)\) commutes with the finite maximum: for any finite nonempty set of real numbers \(\{a_r\}\),
\[
\frac{\max_r a_r}{m(x)}
=
\max_r \frac{a_r}{m(x)}.
\]
Applying this to the displayed maximum, with the zero candidate included as \(0/m(x)=0\), gives
\[
\frac{B_L(x)}{m(x)}
=
\max\left\{
0,\,
\max_{t\in\mathcal Y}
\frac{\ell_{\le t}(x)-h_{\le t}(x)}{m(x)}
\right\}.
\]
% lean: finset_sup'_div_pos

\item
Similarly, \cref{obj:def:threshold-cuts} gives
\[
B_U(x)
=
\min\left\{
m(x),\
\min_{t\in\mathcal Y}\bigl[\ell_{<t}(x)+h_{>t}(x)\bigr]
\right\}.
\]
Division by \(m(x)>0\) commutes with the finite minimum:
\[
\frac{\min_r a_r}{m(x)}
=
\min_r \frac{a_r}{m(x)}.
\]
The cap contributes \(m(x)/m(x)=1\), and the remaining candidates scale termwise.  Therefore
\[
\frac{B_U(x)}{m(x)}
=
\min\left\{
1,\,
\min_{t\in\mathcal Y}
\frac{\ell_{<t}(x)+h_{>t}(x)}{m(x)}
\right\}.
\]
These two normalized expressions are the fixed-marginal ordinal strict-benefit formulas recorded by \citep[Proposition 2, equation (6), p.~546]{LuDingDasgupta2018}.
% lean: finset_inf'_div_pos
\end{enumerate}
\end{proof}

\begin{proof}[Proof of \cref{obj:thm:sharp-exact-mass-threshold-interval}]
Let
\[
c=\{\ell_i(x),h_j(x),q_0(x),q_1(x),\Delta q(x),m(x)\}_{x,i,j}
\]
be the observable capacity vector associated with \(P_{\mathrm{obs}}\).  By
\cref{obj:prop:capacity-identification}, \(c\) is a valid capacity vector,
\(p_x\ge 0\) for every \(x\), and the positive aggregate survivor condition gives
\[
M=\sum_{x\in\mathcal X}p_xm(x)>0 .
\]
For a coupling \(\gamma_x\), write
\[
b_x(\gamma_x)=\sum_{i\in\mathcal Y}\sum_{j\in\mathcal Y:\,i<j}\gamma_{ij,x}.
\]

\begin{enumerate}
\item Fix a cell \(x\) with \(p_x>0\).  If \(W\) is any feasible full law with observed law \(P_{\mathrm{obs}}\), its survivor-complier coupling
\[
\gamma_{ij,x}
=
P_W(Y_0=i,Y_1=j,S_0=S_1=1,C\mid X=x)
\]
is nonnegative.  Its row sums are bounded by the lower-arm selected-complier capacities, its column sums are bounded by the upper-arm selected-complier capacities, and its total mass is the exact survivor-complier mass:
\[
\sum_j\gamma_{ij,x}\le \ell_i(x),\qquad
\sum_i\gamma_{ij,x}\le h_j(x),\qquad
\sum_{i,j}\gamma_{ij,x}=m(x).
\]
Thus every feasible full law projects into \(\Gamma_x^\ast\).

Conversely, take any \(\gamma_x\in\Gamma_x^\ast\).  Define a cellwise family of couplings by
\[
\widetilde\gamma_{x'}=
\begin{cases}
\gamma_x, & x'=x,\\
\gamma_{x'}^L, & x'\ne x,
\end{cases}
\]
where \(\gamma_{x'}^L\) is the lower threshold-flow coupling from
\cref{obj:def:threshold-flow-construction}.  By
\cref{obj:thm:linear-sparse-threshold-flow},
\[
\gamma_{x'}^L\in\Gamma_{x'}^\ast
\]
for every \(x'\), and the chosen \(\gamma_x\) lies in \(\Gamma_x^\ast\) by hypothesis.  Hence
\[
\widetilde\gamma_{x'}\in\Gamma_{x'}^\ast\qquad\text{for every }x'\in\mathcal X .
\]
We record the realization step as a proof-local realization claim.  Let \(\bar\gamma=\{\bar\gamma_{x'}:x'\in\mathcal X\}\) be any family satisfying
\[
\bar\gamma_{x'}\in\Gamma_{x'}^\ast\qquad\text{for every }x'\in\mathcal X .
\]
Choose one compatible generating law \(\bar P\) for \(P_{\mathrm{obs}}\).  For each supported cell \(x'\), put
\[
\mu_C(x')=\bar P(C\mid X=x')
\]
and choose a fixed outcome level \(y_\star\in\mathcal Y\).  Define the one-sided residual masses
\[
\rho_i^0(x')=
\begin{cases}
\ell_i(x')-\sum_j\bar\gamma_{ij,x'}, & \Delta q(x')<0,\\
0, & \Delta q(x')\ge0,
\end{cases}
\]
\[
\rho_j^1(x')=
\begin{cases}
h_j(x')-\sum_i\bar\gamma_{ij,x'}, & \Delta q(x')>0,\\
0, & \Delta q(x')\le0,
\end{cases}
\]
and
\[
\rho^\varnothing(x')=\mu_C(x')-\max\{q_0(x'),q_1(x')\}.
\]
These quantities are nonnegative on every supported cell.  The nonnegativity of \(\rho_i^0\) and \(\rho_j^1\) follows from the row and column inequalities defining \(\Gamma_{x'}^\ast\).  For \(\rho^\varnothing\), \cref{obj:prop:capacity-identification} and \cref{obj:def:observable-capacities} give
\[
q_0(x')=\sum_i\ell_i(x')
=\sum_i \bar P(Y_0=i,S_0=1,C\mid X=x')
=\bar P(S_0=1,C\mid X=x')
\le \bar P(C\mid X=x')=\mu_C(x')
\]
and similarly
\[
q_1(x')=\sum_j h_j(x')
=\sum_j \bar P(Y_1=j,S_1=1,C\mid X=x')
=\bar P(S_1=1,C\mid X=x')
\le \bar P(C\mid X=x')=\mu_C(x').
\]
The equalities summing over outcomes use the finite ordered outcome support, and the inequalities use inclusion of each selected-complier event in the complier event.  Hence \(\max\{q_0(x'),q_1(x')\}\le\mu_C(x')\).  The residual formulas are used on supported cells; cells with \(p_{x'}=0\) are assigned a zero-weight normalized table below.

The residuals complete \(\bar\gamma_{x'}\) to the selected-complier margins.  If \(\Delta q(x')\ge0\), then \(m(x')=q_0(x')\), and the row inequalities in \(\Gamma_{x'}^\ast\) have total slack zero; if \(\Delta q(x')<0\), the definition of \(\rho_i^0\) fills the row slack.  Thus, in all cases,
\[
\sum_j\bar\gamma_{ij,x'}+\rho_i^0(x')=\ell_i(x')\qquad\text{for every }i.
\]
Similarly, if \(\Delta q(x')\le0\), then \(m(x')=q_1(x')\), and the column inequalities have total slack zero; if \(\Delta q(x')>0\), the definition of \(\rho_j^1\) fills the column slack.  Hence
\[
\sum_i\bar\gamma_{ij,x'}+\rho_j^1(x')=h_j(x')\qquad\text{for every }j.
\]
The complier masses also add up correctly:
\[
\sum_{i,j}\bar\gamma_{ij,x'}+
\sum_i\rho_i^0(x')+
\sum_j\rho_j^1(x')+\rho^\varnothing(x')
=\mu_C(x').
\]
Indeed, the left side equals \(m(x')+q_0(x')-m(x')+\mu_C(x')-q_0(x')\) when \(\Delta q(x')<0\), equals \(m(x')+q_1(x')-m(x')+\mu_C(x')-q_1(x')\) when \(\Delta q(x')>0\), and equals \(m(x')+\mu_C(x')-m(x')\) when \(\Delta q(x')=0\).

Construct a latent table in cell \(x'\) by assigning \(\bar\gamma_{ij,x'}\) to compliers with \((S_0,S_1)=(1,1)\) and \((Y_0,Y_1)=(i,j)\), assigning \(\rho_i^0(x')\) to compliers with \((S_0,S_1)=(1,0)\), \(Y_0=i\), and \(Y_1=y_\star\), assigning \(\rho_j^1(x')\) to compliers with \((S_0,S_1)=(0,1)\), \(Y_0=y_\star\), and \(Y_1=j\), and assigning \(\rho^\varnothing(x')\) to compliers with \((S_0,S_1)=(0,0)\) and \((Y_0,Y_1)=(y_\star,y_\star)\).  Keep the never-taker and always-taker conditional latent components from \(\bar P\).  On cells with \(p_{x'}=0\), choose any normalized latent table, since those cells have zero joint weight.

Combining these conditional tables with the original cell weights \(p_{x'}\) and instrument propensities gives a full law \(W^{\bar\gamma}\).  The preceding displays show that the selected-complier outcome margins are \(\ell_i(x')\) and \(h_j(x')\), the complier total is \(\mu_C(x')\), and the noncomplier components are unchanged; therefore the observed law is \(P_{\mathrm{obs}}\).  The construction makes the instrument conditionally independent given \(X\), and the observed variables are generated from the potential variables by the consistency and exclusion equations.  Treatment monotonicity is inherited from the retained principal strata.  The one-sided residuals respect the weak-selection direction because \cref{obj:prop:capacity-identification} gives the sign implications for \(\Delta q(x')\).  Thus \(W^{\bar\gamma}\) is feasible, and for every supported cell
\[
P_{W^{\bar\gamma}}(Y_0=i,Y_1=j,S_0=S_1=1,C\mid X=x')=\bar\gamma_{ij,x'}.
\]
Moreover, since \(\sum_{i,j}\bar\gamma_{ij,x'}=m(x')\), the same cell decomposition gives
\[
\Pr_{W^{\bar\gamma}}(Y_1>Y_0\mid S_0=S_1=1,C)
=
\frac{\sum_{x'}p_{x'}b_{x'}(\bar\gamma_{x'})}{M}.
\]


Applying this proof-local realization claim to \(\widetilde\gamma\), the realized coupling in cell \(x\) is \(\gamma_x\).  Therefore the projection set of feasible full laws is exactly \(\Gamma_x^\ast\). % lean: sharp_exact_mass_threshold_interval

\item Again fix \(x\) with \(p_x>0\).  If \(\Delta q(x)>0\), then \(q_0(x)\le q_1(x)\).  For any \(\gamma_x\in\Gamma_x^\ast\),
\[
\sum_i\sum_j\gamma_{ij,x}=m(x)=q_0(x)=\sum_i\ell_i(x),
\]
while each row satisfies \(\sum_j\gamma_{ij,x}\le \ell_i(x)\).  Since all row deficits are nonnegative and their sum is zero,
\[
\sum_j\gamma_{ij,x}=\ell_i(x)\qquad\text{for every }i.
\]
Thus \(\Gamma_x^\ast\subseteq\Gamma_x^+\), and the reverse inclusion follows because exact rows imply total mass \(q_0(x)=m(x)\) while retaining nonnegativity and the column bounds.  Hence
\[
\Gamma_x^\ast=\Gamma_x^+ .
\]

If \(\Delta q(x)<0\), then \(q_1(x)\le q_0(x)\).  The same argument with columns gives
\[
\sum_i\gamma_{ij,x}=h_j(x)\qquad\text{for every }j
\]
for every \(\gamma_x\in\Gamma_x^\ast\), and exact columns conversely give total mass \(q_1(x)=m(x)\).  Hence
\[
\Gamma_x^\ast=\Gamma_x^- .
\]

If \(\Delta q(x)=0\), then \(q_0(x)=q_1(x)=m(x)\), so both the row and column deficit sums vanish.  Consequently
\[
\sum_j\gamma_{ij,x}=\ell_i(x)\quad\text{and}\quad
\sum_i\gamma_{ij,x}=h_j(x)
\]
for every \(\gamma_x\in\Gamma_x^\ast\), and the converse is immediate from the same total-mass identity.  Therefore
\[
\Gamma_x^\ast=\Gamma_x^0 .
\] % lean: sharp_exact_mass_threshold_interval

\item For every \(\gamma_x\in\Gamma_x^\ast\), the threshold-cut inequalities are
\[
B_L(x)\le b_x(\gamma_x)\le B_U(x).
\]
Indeed, the upper endpoint inequality follows from
\[
b_x(\gamma_x)
\le m(x)
\]
and the threshold cut cover.  Fix \(t\).  For every benefit pair \(i<j\), at least one of \(i<t\) or \(t<j\) holds; otherwise \(t\le i<j\le t\), a contradiction.  Since \(\gamma_x\) is entrywise nonnegative,
\[
b_x(\gamma_x)
=
\sum_i\sum_{j:\,i<j}\gamma_{ij,x}
\le
\sum_{i<t}\sum_j\gamma_{ij,x}
+
\sum_{j>t}\sum_i\gamma_{ij,x}.
\]
The row and column bounds in \(\Gamma_x^\ast\) then give
\[
\sum_{i<t}\sum_j\gamma_{ij,x}
+
\sum_{j>t}\sum_i\gamma_{ij,x}
\le
\ell_{<t}(x)+h_{>t}(x).
\]
Together with the bound by \(m(x)\), this gives \(b_x(\gamma_x)\le B_U(x)\).  For the lower bound, nonnegativity gives \(0\le b_x(\gamma_x)\).  Fix \(t\).  Entry by entry,
\[
\bigl(\mathbf 1\{i\le t\}-\mathbf 1\{j\le t\}\bigr)\gamma_{ij,x}
\le
\mathbf 1\{i<j\}\gamma_{ij,x},
\]
where \(\mathbf 1\{\cdot\}\) denotes an indicator.  If \(i\le t<j\), the two sides are equal because \(i<j\); if \(i\le t\) and \(j\le t\), the left side is zero; if \(t<i\) and \(j\le t\), the left side is nonpositive; and if \(t<i\) and \(t<j\), the left side is zero.  Nonnegativity of \(\gamma_{ij,x}\) handles the nonpositive and zero cases.  Summing this entrywise comparison gives
\[
\sum_{i\le t}\sum_j\gamma_{ij,x}
-
\sum_{j\le t}\sum_i\gamma_{ij,x}
\le b_x(\gamma_x).
\]
When \(\Delta q(x)\ge0\), the row sums bind by the preceding step and
\(\sum_{j\le t}\sum_i\gamma_{ij,x}\le h_{\le t}(x)\), so
\[
\ell_{\le t}(x)-h_{\le t}(x)\le b_x(\gamma_x).
\]
When \(\Delta q(x)<0\), the column sums bind and the total row slack is
\[
\sum_i\left(\ell_i(x)-\sum_j\gamma_{ij,x}\right)=q_0(x)-m(x)=-\Delta q(x).
\]
The prefix row slack is at most this total slack, and therefore
\[
\ell_{\le t}(x)-h_{\le t}(x)+\Delta q(x)\le b_x(\gamma_x).
\]
Combining the two cases gives
\[
\ell_{\le t}(x)-h_{\le t}(x)+\min\{\Delta q(x),0\}\le b_x(\gamma_x)
\]
for every \(t\), and hence \(B_L(x)\le b_x(\gamma_x)\).

For the lower and upper threshold-flow couplings from
\cref{obj:def:threshold-flow-construction}, \cref{obj:thm:linear-sparse-threshold-flow} gives
\[
\gamma_x^L\in\Gamma_x^\ast,\qquad b_x(\gamma_x^L)=B_L(x),
\qquad
\gamma_x^U\in\Gamma_x^\ast,\qquad b_x(\gamma_x^U)=B_U(x).
\]
Thus
\[
\inf_{\gamma_x\in\Gamma_x^\ast}b_x(\gamma_x)=B_L(x),
\qquad
\sup_{\gamma_x\in\Gamma_x^\ast}b_x(\gamma_x)=B_U(x).
\] % lean: sharp_exact_mass_threshold_interval

\item If \(m(x)=0\), then any \(\gamma_x\in\Gamma_x^\ast\) has
\[
\sum_{i,j}\gamma_{ij,x}=0
\]
and all entries are nonnegative.  Hence every entry is zero, so \(\Gamma_x^\ast\subseteq\{0\}\).  Conversely, the zero coupling is nonnegative, has all row and column sums equal to zero, and has total mass zero.  Since the capacity vector is valid, \(\ell_i(x)\ge0\) and \(h_j(x)\ge0\) for every \(i,j\), so the row and column bounds hold; the total-mass constraint holds because \(m(x)=0\).  Thus \(0\in\Gamma_x^\ast\), and therefore
\[
\Gamma_x^\ast=\{0\}.
\] % lean: sharp_exact_mass_threshold_interval

\item Define
\[
\theta_L=\frac{\sum_xp_xB_L(x)}{M},
\qquad
\theta_U=\frac{\sum_xp_xB_U(x)}{M}.
\]
The cellwise endpoint inequality \(B_L(x)\le B_U(x)\) follows by applying the preceding bounds to \(\gamma_x^L\), and \(p_x\ge0\) gives
\[
\sum_xp_xB_L(x)\le \sum_xp_xB_U(x),
\qquad
\theta_L\le\theta_U .
\]

Let \(W\) be any feasible full law with observed law \(P_{\mathrm{obs}}\).  Equality of observed laws transfers the cell weights:
\[
p_x^W=p_x .
\]
For supported cells, the projected coupling \(\gamma_x^W\) lies in \(\Gamma_x^\ast\); for unsupported cells, \(p_x=0\).  Therefore
\[
\sum_xp_xB_L(x)
\le
\sum_xp_x^W b_x(\gamma_x^W)
\le
\sum_xp_xB_U(x),
\]
and the denominator of the benefit probability is
\[
\sum_xp_x^W\sum_{i,j}\gamma_{ij,x}^W
=
\sum_xp_xm(x)
=
M .
\]
The aggregate identity for the benefit probability gives
\[
\Pr_W(Y_1>Y_0\mid S_0=S_1=1,C)
=
\frac{\sum_xp_x^W b_x(\gamma_x^W)}{M},
\]
so every feasible value belongs to \([\theta_L,\theta_U]\).

For the reverse inclusion, take a real number \(\vartheta\in[\theta_L,\theta_U]\).  If \(\theta_L=\theta_U\), choose the family \(\bar\gamma_x=\gamma_x^L\).  The proof-local realization lemma gives a feasible law \(W^L\) with
\[
\Pr_{W^L}(Y_1>Y_0\mid S_0=S_1=1,C)
=
\frac{\sum_xp_xb_x(\gamma_x^L)}{M}
=
\theta_L
=
\vartheta .
\]
If \(\theta_L<\theta_U\), set
\[
\lambda=\frac{\vartheta-\theta_L}{\theta_U-\theta_L},
\qquad 0\le\lambda\le1,
\]
and define
\[
\gamma_x^\lambda=(1-\lambda)\gamma_x^L+\lambda\gamma_x^U .
\]
Since nonnegativity, the row bounds, the column bounds, and the exact total-mass identity are preserved by convex combinations,
\[
\gamma_x^\lambda\in\Gamma_x^\ast\qquad\text{for every }x.
\]
Linearity of \(b_x\) gives
\[
b_x(\gamma_x^\lambda)
=
(1-\lambda)B_L(x)+\lambda B_U(x).
\]
The proof-local realization lemma therefore yields a feasible law \(W^\lambda\) with
\[
\Pr_{W^\lambda}(Y_1>Y_0\mid S_0=S_1=1,C)
=
\frac{\sum_xp_x\{(1-\lambda)B_L(x)+\lambda B_U(x)\}}{M}
=
(1-\lambda)\theta_L+\lambda\theta_U
=
\vartheta .
\]
Hence the feasible values are exactly \([\theta_L,\theta_U]\), which is the interval in \cref{obj:def:identified-interval}. % lean: sharp_exact_mass_threshold_interval
\end{enumerate}
\end{proof}

\begin{proof}[Proof of \cref{obj:prop:three-level-witness}]
Let \(c\) be the observable capacity vector obtained from \(P_{\mathrm{obs}}^{\star}\) on the singleton covariate space \(\{x^{\star}\}\) and outcome levels \(\{0,1,2\}\).

\begin{enumerate}
\item Consider the latent allocation from \cref{obj:def:three-level-witness}, written with the canonical hidden-outcome completion
\[
\begin{array}{c|c|c|c}
(D_0,D_1) & (S_0,S_1) & (Y_0,Y_1) & \text{mass}\\ \hline
(0,1)&(1,1)&(0,0)&3/40\\
(0,1)&(1,1)&(1,1)&1/10\\
(0,1)&(1,1)&(2,2)&3/40\\
(0,1)&(0,1)&(0,0)&1/40\\
(0,1)&(0,1)&(0,1)&3/20\\
(0,1)&(0,1)&(0,2)&3/40\\
(0,1)&(0,0)&(0,0)&1/2 .
\end{array}
\]
The masses are nonnegative and sum to
\[
\frac{3}{40}+\frac{1}{10}+\frac{3}{40}
+\frac{1}{40}+\frac{3}{20}+\frac{3}{40}
+\frac12
=
1.
\]
Let \(Z\) be independent of these potential variables with
\(P(Z=0)=P(Z=1)=1/2\), and set \(D=D_Z\), \(S=S_D\), and \(Y=Y_D\) on selected units.  Every latent atom has \(D_0=0\), \(D_1=1\), and \(S_0\le S_1\).  The instrument propensity in the only cell is \(1/2\), so the overlap inequalities at \(\varepsilon_Z=1/4\) are
\[
0<\frac14<\frac12,
\qquad
\frac14\le \frac12\le \frac34 .
\]
The always-selected complier mass is
\[
\frac{3}{40}+\frac{1}{10}+\frac{3}{40}=\frac14>0.
\]
Thus the law satisfies the structural requirements collected in \cref{obj:def:structural-law-class}, with the single cell retained and weak selection direction \(S_0\le S_1\).

Under \(Z=0\), the selected outcome masses are \(3/40,1/10,3/40\), and the unselected mass is \(3/4\).  Multiplying by \(P(Z=0)=1/2\) gives the first four observed masses
\[
\frac{3}{80},\quad \frac{1}{20},\quad \frac{3}{80},\quad \frac38 .
\]
Under \(Z=1\), the selected outcome masses are
\[
\left(\frac{3}{40},\frac{1}{10},\frac{3}{40}\right)
+
\left(\frac{1}{40},\frac{3}{20},\frac{3}{40}\right)
=
\left(\frac{1}{10},\frac14,\frac{3}{20}\right),
\]
and the unselected mass is \(1/2\).  Multiplying by \(P(Z=1)=1/2\) gives the last four observed masses
\[
\frac{1}{20},\quad \frac18,\quad \frac{3}{40},\quad \frac14 .
\]
Hence this three-level slate system realizes \(P_{\mathrm{obs}}^{\star}\), has the displayed encoded full law, and satisfies \(D_0=0\), \(D_1=1\), and \(S_0\le S_1\) almost surely. % lean: wRealization

\item By \cref{obj:def:observable-capacities}, for \(i\in\{0,1,2\}\),
\[
\ell_i(x^{\star})
=
P(D=0,S=1,Y=i\mid Z=0,X=x^{\star})
-
P(D=0,S=1,Y=i\mid Z=1,X=x^{\star}).
\]
The second conditional probability is \(0\) for each \(i\), and the \(Z=0\) denominator is \(1/2\).  Therefore
\[
(\ell_0(x^{\star}),\ell_1(x^{\star}),\ell_2(x^{\star}))
=
\left(
\frac{(3/80)}{1/2},
\frac{(1/20)}{1/2},
\frac{(3/80)}{1/2}
\right)
=
\left(\frac{3}{40},\frac{1}{10},\frac{3}{40}\right).
\]
% lean: wLower

\item Similarly, for \(j\in\{0,1,2\}\),
\[
h_j(x^{\star})
=
P(D=1,S=1,Y=j\mid Z=1,X=x^{\star})
-
P(D=1,S=1,Y=j\mid Z=0,X=x^{\star}).
\]
The second conditional probability is \(0\) for each \(j\), and the \(Z=1\) denominator is \(1/2\).  Hence
\[
(h_0(x^{\star}),h_1(x^{\star}),h_2(x^{\star}))
=
\left(
\frac{(1/20)}{1/2},
\frac{(1/8)}{1/2},
\frac{(3/40)}{1/2}
\right)
=
\left(\frac{1}{10},\frac14,\frac{3}{20}\right).
\]
% lean: wUpper

\item Summing the displayed capacities gives
\[
q_0(x^{\star})=\frac{3}{40}+\frac{1}{10}+\frac{3}{40}=\frac14,
\qquad
q_1(x^{\star})=\frac{1}{10}+\frac14+\frac{3}{20}=\frac12.
\]
Thus
\[
m(x^{\star})=\min\{q_0(x^{\star}),q_1(x^{\star})\}=\frac14,
\qquad
\Delta q(x^{\star})=q_1(x^{\star})-q_0(x^{\star})=\frac14 .
\]
Using the threshold-cut formulas in \cref{obj:def:threshold-cuts}, the lower-cut correction is
\[
\min\{\Delta q(x^{\star}),0\}=0.
\]
The three prefix differences are
\[
\ell_{\le0}(x^{\star})-h_{\le0}(x^{\star})=-\frac{1}{40},
\]
\[
\ell_{\le1}(x^{\star})-h_{\le1}(x^{\star})=-\frac{7}{40},
\]
\[
\ell_{\le2}(x^{\star})-h_{\le2}(x^{\star})=-\frac14 .
\]
The maximum with \(0\) is therefore
\[
B_L(x^{\star})=0.
\]
For the upper cut,
\[
m(x^{\star})=\frac{10}{40},
\]
and the strict-prefix/tail sums are
\[
\ell_{<0}(x^{\star})+h_{>0}(x^{\star})=\frac25,
\]
\[
\ell_{<1}(x^{\star})+h_{>1}(x^{\star})=\frac{9}{40},
\]
\[
\ell_{<2}(x^{\star})+h_{>2}(x^{\star})=\frac{7}{40}.
\]
Taking the minimum gives
\[
B_U(x^{\star})
=
\min\left\{\frac{10}{40},\frac25,\frac{9}{40},\frac{7}{40}\right\}
=
\frac{7}{40}.
\]
% lean: wCapacityFacts

\item The capacity vector is valid because each displayed coordinate of \(\ell\) and \(h\) is nonnegative.  The singleton cell has observed weight
\[
p_{x^{\star}}=1\ge0,
\]
and the aggregate survivor-complier mass is
\[
M=\sum_x p_xm(x)=1\cdot \frac14=\frac14>0.
\]
Since the observed law is induced by the latent slate system constructed above, it lies in the observed-law domain.  By \cref{obj:def:identified-interval},
\[
\Theta_I(P_{\mathrm{obs}}^{\star})
=
\left[
\frac{p_{x^{\star}}B_L(x^{\star})}{M},
\frac{p_{x^{\star}}B_U(x^{\star})}{M}
\right]
=
\left[
\frac{0}{1/4},
\frac{(7/40)}{1/4}
\right]
=
\left[0,\frac{7}{10}\right].
\]
% lean: wInterval

\item The hypotheses of \cref{obj:thm:full-law-endpoint-attainment} hold for the realized three-level slate system: the structural restrictions and tie-safe survivor model hold by the first step, and \(M=1/4>0\) by the preceding step.  Applying that result to \(P_{\mathrm{obs}}^{\star}\), overlap \(1/4\), and the retained singleton cell gives full laws \(P^L\) and \(P^U\), with three-level slate systems, that induce \(P_{\mathrm{obs}}^{\star}\) and attain the two endpoint values
\[
0
\qquad\text{and}\qquad
\frac{7}{10}.
\]
% lean: wEndpointWitnesses
\end{enumerate}
Combining the six displayed conclusions gives all assertions of \cref{obj:prop:three-level-witness}. % lean: three_level_witness_sharp
\end{proof}

\begin{proof}[Proof of \cref{obj:thm:linear-sparse-threshold-flow}]
Choose
\[
C_{\mathrm{sparse}}=75,
\qquad
C_{\mathrm{dense}}=77 .
\]
Both constants are positive.  Fix a finite covariate support \(\mathcal X\), an integer \(K\ge 3\), and nonnegative capacity arrays \(\ell_i(x)\) and \(h_j(x)\).  Write \(N_{\mathcal X}=|\mathcal X|\).  For one cell \(x\), let
\[
s_K=64K+32
\]
be the fixed sparse cell budget charged by the implementation.  Since \(K\ge3\),
\[
s_K=64K+32\le 64K+11K=75K .
\]
% lean: linear_sparse_threshold_flow

\begin{enumerate}
\item The actual sparse work in cell \(x\) is the sum of the threshold-cut scan, the lower nested-graph pass, the upper nested-graph pass, and the two completion passes.  The threshold-cut scan uses
\[
12K+8
\]
operations.  If \(\mathsf{Rows}\) and \(\mathsf{Cols}\) are the current row and column lists, each nested pass performs at most
\[
5\bigl(|\mathsf{Rows}|+|\mathsf{Cols}|\bigr)+1
\]
operations, and each completion pass performs at most
\[
4\bigl(|\mathsf{Rows}|+|\mathsf{Cols}|\bigr)+1
\]
operations.  Initially the row and column lists have lengths \(K\) and \(K\).  After either nested pass, the residual row and column lists have total length at most \(2K\).  Hence, for the actual cell count \(a_x\),
\[
\begin{aligned}
a_x
&\le (12K+8)+\{5(2K)+1\}+\{5(2K)+1\}
      +\{4(2K)+1\}+\{4(2K)+1\}  \\
&=48K+12
\le 64K+32=s_K .
\end{aligned}
\]
Consequently
\[
\sum_{x\in\mathcal X} a_x
\le \sum_{x\in\mathcal X}s_K
\le \sum_{x\in\mathcal X}75K
=75\,N_{\mathcal X}K .
\]
This proves the actual sparse bound.  The charged sparse count is exactly
\[
\sum_{x\in\mathcal X}s_K,
\]
so the same calculation gives
\[
\sum_{x\in\mathcal X}s_K\le 75\,N_{\mathcal X}K .
\]
% lean: sparseCellActualOperations_le_budget

\item Since \(K\ge3\), in particular \(1\le K\), and therefore
\[
K\le K^2 .
\]
Multiplying the sparse bounds by the nonnegative factor \(75N_{\mathcal X}\) gives
\[
75\,N_{\mathcal X}K\le 75\,N_{\mathcal X}K^2 .
\]
The dense actual implementation first performs the sparse computation and then materializes the two dense \(K\times K\) cell matrices, so its actual count is
\[
\sum_{x\in\mathcal X}a_x+2N_{\mathcal X}K^2 .
\]
Using the preceding sparse actual bound,
\[
\sum_{x\in\mathcal X}a_x+2N_{\mathcal X}K^2
\le 75N_{\mathcal X}K^2+2N_{\mathcal X}K^2
=77N_{\mathcal X}K^2 .
\]
The charged dense count is
\[
\sum_{x\in\mathcal X}s_K+2N_{\mathcal X}K^2,
\]
and the same argument gives
\[
\sum_{x\in\mathcal X}s_K+2N_{\mathcal X}K^2
\le 77N_{\mathcal X}K^2 .
\]
% lean: linear_sparse_threshold_flow

\item The charged sparse count is
\[
\sum_{x\in\mathcal X}s_K=\sum_{x\in\mathcal X}(64K+32),
\]
which depends only on \(|\mathcal X|\) and \(K\).  For any other nonnegative capacity arrays on the same \(\mathcal X\) and \(K\), the charged sparse count is the same sum.  The charged dense count adds the same deterministic materialization term:
\[
\sum_{x\in\mathcal X}s_K+2|\mathcal X|K^2 .
\]
Thus the charged sparse counts are equal across such capacity arrays, and the charged dense counts are equal across such capacity arrays.
% lean: linear_sparse_threshold_flow

\item Fix a cell \(x\).  Put \(\mathcal Y\) for the \(K\)-point ordered support indexing the row and column capacities in the statement; all sums, extrema, and inequalities among indices below use this finite ordered set.  With these arrays, \cref{obj:def:observable-capacities} gives
\[
q_0(x)=\sum_{i\in\mathcal Y}\ell_i(x),
\qquad
q_1(x)=\sum_{j\in\mathcal Y}h_j(x),
\]
\[
\Delta q(x)=q_1(x)-q_0(x),
\qquad
m(x)=\min\{q_0(x),q_1(x)\}.
\]
Denote the deterministic lower and upper scan outputs by \(B_L^{\mathrm{scan}}(x)\) and \(B_U^{\mathrm{scan}}(x)\), defined by
\[
B_L^{\mathrm{scan}}(x)
=
\max\left\{0,\max_{t\in\mathcal Y}
\left[\ell_{\le t}(x)-h_{\le t}(x)+\min\{\Delta q(x),0\}\right]\right\}
\]
and
\[
B_U^{\mathrm{scan}}(x)
=
\min\left\{m(x),\min_{t\in\mathcal Y}
\left[\ell_{<t}(x)+h_{>t}(x)\right]\right\}.
\]
The costed sparse routine stores as its value
\[
\bigl((B_L^{\mathrm{scan}}(x),B_U^{\mathrm{scan}}(x)),
       (\tau_x^L,\tau_x^U)\bigr),
\]
where \((\tau_x^L,\tau_x^U)\) are the lower and upper sparse allocation traces.  By \cref{obj:def:threshold-cuts}, these scan outputs are the threshold-cut values:
\[
(B_L^{\mathrm{scan}}(x),B_U^{\mathrm{scan}}(x))=(B_L(x),B_U(x)).
\]
Moreover
\[
\tau_x^L=\text{the sparse lower allocation trace},
\qquad
\tau_x^U=\text{the sparse upper allocation trace}.
\]
% lean: costedThresholdCuts_value

\item Let \(\gamma_x^L\) and \(\gamma_x^U\) be the matrices obtained by materializing the sparse traces \(\tau_x^L\) and \(\tau_x^U\): for \(\bullet\in\{L,U\}\),
\[
\gamma_x^\bullet(i,j)
=
\sum_{e\in\tau_x^\bullet:\ \operatorname{row}(e)=i,\ \operatorname{col}(e)=j}
\operatorname{mass}(e).
\]
For any cellwise coupling \(\gamma\), the benefit-mass functional used below is the strict-benefit sum
\[
b_x(\gamma)
=
\sum_{i\in\mathcal Y}\sum_{\substack{j\in\mathcal Y\\ i<j}}\gamma(i,j).
\]
For either the lower nonbenefit primary pass or the upper benefit primary pass, the input row and column lists have lengths \(K\) and \(K\).  If \(\mathsf{prim}_x^\bullet\) is the primary allocation list and \(\mathsf{row}_{x}^{\bullet,\mathrm{res}}\), \(\mathsf{col}_{x}^{\bullet,\mathrm{res}}\) are the two residual lists, the recursive pass gives
\[
|\mathsf{prim}_x^\bullet|
+|\mathsf{row}_{x}^{\bullet,\mathrm{res}}|
+|\mathsf{col}_{x}^{\bullet,\mathrm{res}}|
\le 2K,
\qquad
0<|\mathsf{row}_{x}^{\bullet,\mathrm{res}}|
+|\mathsf{col}_{x}^{\bullet,\mathrm{res}}|.
\]
The complete bipartite residual pass consumes one residual row or one residual column at each emitted allocation and stops when one residual side is exhausted, so
\[
|\operatorname{complete}(\mathsf{row}_{x}^{\bullet,\mathrm{res}},
\mathsf{col}_{x}^{\bullet,\mathrm{res}})|
<
|\mathsf{row}_{x}^{\bullet,\mathrm{res}}|
+|\mathsf{col}_{x}^{\bullet,\mathrm{res}}|.
\]
Thus
\[
|\tau_x^\bullet|
=
|\mathsf{prim}_x^\bullet|
+|\operatorname{complete}(\mathsf{row}_{x}^{\bullet,\mathrm{res}},
\mathsf{col}_{x}^{\bullet,\mathrm{res}})|
<2K,
\]
and, since \(|\tau_x^\bullet|\) is an integer,
\[
|\tau_x^L|\le 2K-1,
\qquad
|\tau_x^U|\le 2K-1.
\]
Materializing a sparse trace can only put positive mass on a listed row-column pair, hence
\[
|\{(i,j):\gamma_x^L(i,j)>0\}|\le |\tau_x^L|\le 2K-1,
\qquad
|\{(i,j):\gamma_x^U(i,j)>0\}|\le |\tau_x^U|\le 2K-1 .
\]
% lean: thresholdFlowLower_positiveSupportCard_le
% lean: thresholdFlowUpper_positiveSupportCard_le

Every emitted mass in the two primary passes and in the residual completion pass is the minimum of two nonnegative residual capacities, and each residual update subtracts that minimum.  Therefore
\[
\gamma_x^L(i,j)\ge0,
\qquad
\gamma_x^U(i,j)\ge0
\]
for all \(i,j\).  Let \(\gamma_{x,\mathrm{prim}}^\bullet\) and \(\gamma_{x,\mathrm{comp}}^\bullet\) be the materialized primary and completion parts, so \(\gamma_x^\bullet=\gamma_{x,\mathrm{prim}}^\bullet+\gamma_{x,\mathrm{comp}}^\bullet\).  For the row and column residual lists left by the primary pass, define the residual capacities
\[
\rho_{x,i}^\bullet
=
\sum_{(i',u)\in\mathsf{row}_{x}^{\bullet,\mathrm{res}}:\ i'=i}u,
\qquad
\chi_{x,j}^\bullet
=
\sum_{(j',v)\in\mathsf{col}_{x}^{\bullet,\mathrm{res}}:\ j'=j}v .
\]
The primary pass conserves each original capacity after its residual is included, while the completion pass uses at most the residual capacities:
\[
\sum_j\gamma_{x,\mathrm{prim}}^\bullet(i,j)+\rho_{x,i}^\bullet=\ell_i(x),
\qquad
\sum_j\gamma_{x,\mathrm{comp}}^\bullet(i,j)\le \rho_{x,i}^\bullet,
\]
and
\[
\sum_i\gamma_{x,\mathrm{prim}}^\bullet(i,j)+\chi_{x,j}^\bullet=h_j(x),
\qquad
\sum_i\gamma_{x,\mathrm{comp}}^\bullet(i,j)\le \chi_{x,j}^\bullet.
\]
Consequently
\[
\sum_j\gamma_x^\bullet(i,j)\le\ell_i(x),
\qquad
\sum_i\gamma_x^\bullet(i,j)\le h_j(x).
\]
The completion transports the smaller of the two residual totals, and the primary pass conserves the original row and column totals.  Hence
\[
\sum_{i,j}\gamma_x^\bullet(i,j)
=
\sum_{i,j}\gamma_{x,\mathrm{prim}}^\bullet(i,j)
+\min\left\{\sum_i\rho_{x,i}^\bullet,\sum_j\chi_{x,j}^\bullet\right\}
=\min\{q_0(x),q_1(x)\}=m(x).
\]
By the definition of \(\Gamma_x^\ast\) in \cref{obj:def:partial-transport-polytope}, these displays give
\[
\gamma_x^L\in\Gamma_x^\ast,
\qquad
\gamma_x^U\in\Gamma_x^\ast.
\]

It remains to identify the two benefit masses.  For the lower pass, let \(\nu_x^L\) be the total mass emitted by the primary nonbenefit pass.  Every primary lower allocation satisfies \(j\le i\), and every residual row-column pair left for completion satisfies \(i<j\).  Hence the primary part contributes zero benefit mass and the completion part contributes all of its transported mass to benefit edges.  Since the completed lower trace transports total mass \(m(x)\),
\[
b_x(\gamma_x^L)=m(x)-\nu_x^L.
\]
For the upper pass, let \(\beta_x^U\) be the total mass emitted by the primary benefit pass.  Every primary upper allocation satisfies \(i<j\), and every residual row-column pair left for completion satisfies \(j\le i\).  Thus the residual completion contributes zero benefit mass and
\[
b_x(\gamma_x^U)=\beta_x^U.
\]
Define the finite lower and upper cut families on \(\{\varnothing\}\cup\mathcal Y\) by
\[
A_x^L(\varnothing)=m(x),
\qquad
A_x^L(t)=\ell_{>t}(x)+h_{\le t}(x),
\]
\[
A_x^U(\varnothing)=m(x),
\qquad
A_x^U(t)=\ell_{<t}(x)+h_{>t}(x),
\]
where \(\ell_{>t}(x)=\sum_{i>t}\ell_i(x)\).  For the lower primary pass, the empty cut bound is
\[
\nu_x^L\le q_0(x),
\qquad
\nu_x^L\le q_1(x),
\qquad
\nu_x^L\le m(x)=A_x^L(\varnothing).
\]
For an ordinary threshold \(t\in\mathcal Y\), each nonbenefit primary allocation is counted by the row side \(i>t\) or by the column side \(j\le t\), so
\[
\nu_x^L
\le
\sum_{i>t}\sum_j\gamma_{x,\mathrm{prim}}^L(i,j)
+
\sum_{j\le t}\sum_i\gamma_{x,\mathrm{prim}}^L(i,j).
\]
Using the primary conservation identities and the nonnegative residual capacities,
\[
\sum_{i>t}\sum_j\gamma_{x,\mathrm{prim}}^L(i,j)
\le \ell_{>t}(x),
\qquad
\sum_{j\le t}\sum_i\gamma_{x,\mathrm{prim}}^L(i,j)
\le h_{\le t}(x),
\]
and therefore
\[
\nu_x^L\le A_x^L(t)
\qquad(t\in\mathcal Y).
\]
Thus
\[
\nu_x^L\le \inf_{t\in\{\varnothing\}\cup\mathcal Y}A_x^L(t).
\]
The reverse inequality is obtained from the boundary invariant of the same ascending pass.  The input row and column lists are ordered increasingly.  By induction on the sum of their lengths, the recursive pass maintains
\[
\text{every emitted primary allocation }(i,j,u)\text{ satisfies }j\le i,
\]
\[
\text{every residual row }(i,u)\text{ and residual column }(j,v)\text{ satisfy }i<j,
\]
and, for every residual row \((r,u)\) and emitted allocation \((i,j,v)\),
\[
\neg(r<i\ \text{and}\ j\le r).
\]
The empty-list cases are immediate.  In an eligible step \(j\le i\), the pass emits \((i,j,\min\{a,b\})\) and recurses after reducing or deleting the exhausted endpoint; the reduced lists preserve the increasing order, and any later residual row has index at least the emitted row index, so the new emitted allocation cannot cross a later residual-row boundary.  In an ineligible step \(i<j\), the row \(i\) is moved to the residual list; every column that can remain residual has index at least \(j\), and every later emitted allocation uses such a column, giving both the residual separation and the displayed noncrossing condition for this new residual row.

If the lower primary pass leaves no row residual, row-total conservation gives \(\nu_x^L=q_0(x)\), and the column bound above gives \(\nu_x^L\le q_1(x)\).  Thus, using \cref{obj:def:observable-capacities},
\[
\nu_x^L=\min\{q_0(x),q_1(x)\}=m(x)=A_x^L(\varnothing),
\]
so the infimum is at most \(\nu_x^L\).  If a row residual remains, let \(t_L\) be the largest row index among the lower residual rows.  The finite maximum transfers the preceding invariants to the cut:
\[
i'\le t_L\quad\text{for every }(i',u)\in\mathsf{row}_{x}^{L,\mathrm{res}},
\qquad
t_L<j'\quad\text{for every }(j',v)\in\mathsf{col}_{x}^{L,\mathrm{res}},
\]
and no primary lower allocation with positive mass has both \(t_L<i\) and \(j\le t_L\).  Hence
\[
\sum_{i>t_L}\rho_{x,i}^L=0,
\qquad
\sum_{j\le t_L}\chi_{x,j}^L=0.
\]
Since every primary lower allocation satisfies \(j\le i\) and no such allocation crosses the cut \(t_L\), its mass is counted exactly once by the row side \(i>t_L\) or the column side \(j\le t_L\):
\[
\nu_x^L
=
\sum_{i>t_L}\sum_j\gamma_{x,\mathrm{prim}}^L(i,j)
+
\sum_{j\le t_L}\sum_i\gamma_{x,\mathrm{prim}}^L(i,j).
\]
Adding the zero residual sums and applying row and column conservation gives
\[
\nu_x^L
=
\ell_{>t_L}(x)+h_{\le t_L}(x)
=A_x^L(t_L),
\]
so again \(\inf_{t\in\{\varnothing\}\cup\mathcal Y}A_x^L(t)\le \nu_x^L\).  Consequently
\[
\nu_x^L
=
\inf_{t\in\{\varnothing\}\cup\mathcal Y} A_x^L(t).
\]
For the upper primary pass the cut bounds use the strict-benefit cuts.  The empty cut gives
\[
\beta_x^U\le q_0(x),
\qquad
\beta_x^U\le q_1(x),
\qquad
\beta_x^U\le m(x)=A_x^U(\varnothing),
\]
and, for every \(t\in\mathcal Y\), each benefit primary allocation is counted by the row side \(i<t\) or by the column side \(j>t\):
\[
\beta_x^U
\le
\sum_{i<t}\sum_j\gamma_{x,\mathrm{prim}}^U(i,j)
+
\sum_{j>t}\sum_i\gamma_{x,\mathrm{prim}}^U(i,j)
\le
\ell_{<t}(x)+h_{>t}(x)
=A_x^U(t).
\]
Thus \(\beta_x^U\le\inf_{t\in\{\varnothing\}\cup\mathcal Y}A_x^U(t)\).  For the reverse inequality, use the symmetric boundary invariant for the descending benefit pass.  Its input row and column lists are ordered decreasingly, and induction on the same recursion gives
\[
\text{every emitted primary allocation }(i,j,u)\text{ satisfies }i<j,
\]
\[
\text{every residual row }(i,u)\text{ and residual column }(j,v)\text{ satisfy }j\le i,
\]
and, for every residual row \((r,u)\) and emitted allocation \((i,j,v)\),
\[
\neg(i<r\ \text{and}\ r<j).
\]
The eligible branch \(i<j\) emits \((i,j,\min\{a,b\})\) and recurses on decreasing lists after the exhausted endpoint is reduced or deleted; any later residual row has index at most the emitted row index, so the new emitted allocation cannot cross its residual-row boundary.  In the ineligible branch \(j\le i\), the row \(i\) is moved to the residual list; all remaining columns have index at most \(j\), and all later emitted allocations use such columns, giving the residual separation and noncrossing condition for the new residual row.

If the upper primary pass leaves no row residual, then \(\beta_x^U=q_0(x)\) and the column bound above gives \(\beta_x^U\le q_1(x)\), so, by \cref{obj:def:observable-capacities},
\[
\beta_x^U=\min\{q_0(x),q_1(x)\}=m(x)=A_x^U(\varnothing).
\]
If a row residual remains, let \(t_U\) be the smallest row index among the upper residual rows.  The finite minimum transfers the preceding invariants to the cut:
\[
t_U\le i'\quad\text{for every }(i',u)\in\mathsf{row}_{x}^{U,\mathrm{res}},
\qquad
j'\le t_U\quad\text{for every }(j',v)\in\mathsf{col}_{x}^{U,\mathrm{res}},
\]
and no primary upper allocation with positive mass has both \(i<t_U\) and \(t_U<j\).  Hence
\[
\sum_{i<t_U}\rho_{x,i}^U=0,
\qquad
\sum_{j>t_U}\chi_{x,j}^U=0.
\]
Since every primary upper allocation satisfies \(i<j\) and no such allocation crosses the cut \(t_U\), its mass is counted exactly once by the row side \(i<t_U\) or the column side \(j>t_U\):
\[
\beta_x^U
=
\sum_{i<t_U}\sum_j\gamma_{x,\mathrm{prim}}^U(i,j)
+
\sum_{j>t_U}\sum_i\gamma_{x,\mathrm{prim}}^U(i,j).
\]
Conservation then yields
\[
\beta_x^U
=
\ell_{<t_U}(x)+h_{>t_U}(x)
=A_x^U(t_U).
\]
Therefore
\[
\beta_x^U
=
\inf_{t\in\{\varnothing\}\cup\mathcal Y} A_x^U(t).
\]

For the lower endpoint, the partition identity \(\ell_{\le t}(x)+\ell_{>t}(x)=q_0(x)\), together with \(m(x)=q_0(x)+\min\{\Delta q(x),0\}\) from \cref{obj:def:observable-capacities}, implies
\[
m(x)-A_x^L(\varnothing)=0,
\qquad
m(x)-A_x^L(t)
=\ell_{\le t}(x)-h_{\le t}(x)+\min\{\Delta q(x),0\}.
\]
Therefore, by \cref{obj:def:threshold-cuts},
\[
b_x(\gamma_x^L)
=m(x)-\inf_t A_x^L(t)
=\sup_t\{m(x)-A_x^L(t)\}
=B_L(x).
\]
For the upper endpoint, the displayed formula for \(A_x^U\) is exactly the upper threshold-cut family in \cref{obj:def:threshold-cuts}, and hence
\[
b_x(\gamma_x^U)
=\inf_t A_x^U(t)
=B_U(x).
\]
Since \(x\) was arbitrary, the asserted cellwise output properties hold in every covariate cell.
% lean: thresholdFlow_endpoint_spec
\end{enumerate}
\end{proof}

\begin{proof}[Proof of \cref{obj:thm:branch-free-pointwise-directional-limit}]
\begin{enumerate}
\item Let \(\mathcal O\) be the finite set of observed cells and write
\[
\vartheta_o=P_{\mathrm{obs}}(\{o\}),
\qquad
\widehat\vartheta_{n,o}
=
\frac1n\sum_{r=1}^n\mathbf 1\{O_r=o\},
\qquad o\in\mathcal O .
\]
For an event \(E\subseteq\mathcal O\), set
\[
u(E)=\sum_{o\in E}u_o .
\]
The finite multinomial central limit input applied under \cref{obj:ass:iid-sampling} gives a Gaussian law \(\mathbb Z_{P_{\mathrm{obs}}}\) on \(\mathbb R^{\mathcal O}\) such that
\[
\sqrt n(\widehat\vartheta_n-\vartheta)\rightsquigarrow
\mathbb Z_{P_{\mathrm{obs}}},
\]
\[
\int z\,d\mathbb Z_{P_{\mathrm{obs}}}(z)=0,
\]
and, for observed cells \(a,b\),
\[
\int z(a)z(b)\,d\mathbb Z_{P_{\mathrm{obs}}}(z)
=
\begin{cases}
P_{\mathrm{obs}}(a)\{1-P_{\mathrm{obs}}(a)\},&a=b,\\
-P_{\mathrm{obs}}(a)P_{\mathrm{obs}}(b),&a\ne b.
\end{cases}
\]
Here the displayed empirical vector is the atom vector of the sample empirical law, because its \(o\)-coordinate is exactly the centered empirical frequency of \(\{o\}\). % lean: CausalSmith.PartialID.SlateBenefitPartialTransport.FiniteMultinomialCLT

\item Put
\[
\mathcal X_+(P)=\{x:p_xm(x)>0\}.
\]
For an arbitrary vector \(u\in\mathbb R^{\mathcal O}\), define
\[
\mathsf C_u(A\mid B)=
\begin{cases}
u(A\cap B)/u(B),&u(B)>0,\\
0,&u(B)\le0,
\end{cases}
\]
and form the raw contrasts
\[
\ell_i(u,x)
=
\mathsf C_u(X=x,D=0,S=1,Y=i\mid X=x,Z=0)
-
\mathsf C_u(X=x,D=0,S=1,Y=i\mid X=x,Z=1),
\]
\[
h_j(u,x)
=
\mathsf C_u(X=x,D=1,S=1,Y=j\mid X=x,Z=1)
-
\mathsf C_u(X=x,D=1,S=1,Y=j\mid X=x,Z=0).
\]
Let \(\ell_i^+(u,x)=\max\{\ell_i(u,x),0\}\) and
\(h_j^+(u,x)=\max\{h_j(u,x),0\}\). From these projected capacities define
\(q_0(u,x),q_1(u,x),m(u,x),B_L(u,x),B_U(u,x)\) by the formulas in
\cref{obj:def:observable-capacities,obj:def:threshold-cuts}, and define the fixed-support endpoint map
\[
\Phi(u)
=
\left(
\frac{\sum_{x\in\mathcal X_+(P)}u(X=x)B_L(u,x)}
{\sum_{x\in\mathcal X_+(P)}u(X=x)m(u,x)},
\frac{\sum_{x\in\mathcal X_+(P)}u(X=x)B_U(u,x)}
{\sum_{x\in\mathcal X_+(P)}u(X=x)m(u,x)}
\right),
\]
with the quotient coordinates totalized by the value \(0\) when the common denominator is zero. % lean: CausalSmith.PartialID.SlateBenefitPartialTransport.endpointFromMassVectorOn

\item \cref{obj:prop:capacity-identification} gives compatibility of the observable capacities, \(p_x\ge0\), and the structural identification of \(m(x)\). Since \cref{obj:ass:positive-aggregate-survivors} gives
\[
M=\sum_x p_xm(x)>0,
\]
the recovered support is nonempty after restriction to positive products. If \(x\in\mathcal X_+(P)\), then \(p_xm(x)>0\), hence \(p_x>0\). Writing \(\pi(x)=P_{\mathrm{obs}}(Z=1\mid X=x)\), \cref{obj:ass:instrument-overlap} gives \(0<\varepsilon_Z\), \(\varepsilon_Z\le \pi(x)\), and \(\pi(x)\le 1-\varepsilon_Z\). Therefore
\[
P_{\mathrm{obs}}(X=x,Z=1)=p_x\pi(x)\ge p_x\varepsilon_Z>0
\]
and
\[
P_{\mathrm{obs}}(X=x,Z=0)=p_x\{1-\pi(x)\}\ge p_x\varepsilon_Z>0 .
\]
Thus all conditional denominators used by \(\Phi\) on \(\mathcal X_+(P)\) are positive at \(\vartheta\). Moreover,
\[
\ell_i(\vartheta,x)=\ell_i(x),\qquad h_j(\vartheta,x)=h_j(x),
\]
because \(u(E)\) at \(u=\vartheta\) is \(P_{\mathrm{obs}}(E)\). The nonnegativity supplied by \cref{obj:prop:capacity-identification} then yields
\[
\ell_i^+(\vartheta,x)=\ell_i(x),\qquad h_j^+(\vartheta,x)=h_j(x),
\]
and
\[
\vartheta(X=x)=p_x,\qquad m(\vartheta,x)=m(x).
\]
Capacity validity gives \(\ell_i(x)\ge0\) and \(h_j(x)\ge0\), so
\[
q_0(x)=\sum_i\ell_i(x)\ge0,
\qquad
q_1(x)=\sum_jh_j(x)\ge0,
\qquad
m(x)=\min\{q_0(x),q_1(x)\}\ge0 .
\]
Consequently
\[
\sum_{x\in\mathcal X_+(P)}\vartheta(X=x)m(\vartheta,x)
=
\sum_{x\in\mathcal X_+(P)}p_xm(x)
=
M>0 .
\]
For the last equality, if \(x\notin\mathcal X_+(P)\), then \(\neg(0<p_xm(x))\); since \cref{obj:prop:capacity-identification} gives \(p_x\ge0\) and the preceding display gives \(m(x)\ge0\), we have \(0\le p_xm(x)\), hence \(p_xm(x)=0\). % lean: CausalSmith.PartialID.SlateBenefitPartialTransport.capacity_identification

\item The map \(u\mapsto\Phi(u)\) is Hadamard directionally differentiable at \(\vartheta\) along all directions in \(\mathbb R^{\mathcal O}\). Indeed, on the neighborhood where the arm denominators in the preceding step are positive, the raw conditional contrasts are smooth quotient maps. The coordinatewise positive-part map has scalar directional derivative
\[
\dot a\mapsto
\begin{cases}
\dot a,&a>0,\\
\max\{\dot a,0\},&a=0,
\end{cases}
\]
at every nonnegative coordinate \(a\). Finite sums preserve directional differentiability, the scalar minimum gives the derivative of \(m(u,x)=\min\{q_0(u,x),q_1(u,x)\}\), and the finite maximum and minimum defining \(B_L(u,x)\) and \(B_U(u,x)\) give the active-face directional derivatives of the threshold cuts. Therefore the denominator and the two numerators in \(\Phi\) are continuously Hadamard directionally differentiable at \(\vartheta\). Since the denominator is positive at \(\vartheta\), the quotient rule gives a continuous directional derivative \(D\Psi\) satisfying
\[
\frac{\Phi(\vartheta+t_n g_n)-\Phi(\vartheta)}{t_n}
\longrightarrow
D\Psi(g)
\]
whenever \(t_n\downarrow0\) and \(g_n\to g\) in \(\mathbb R^{\mathcal O}\). Applying the directional delta-method input to \(\Phi\) gives
\[
\sqrt n\{\Phi(\widehat\vartheta_n)-\Phi(\vartheta)\}
\rightsquigarrow
D\Psi(\mathbb Z_{P_{\mathrm{obs}}}).
\]
% lean: CausalSmith.PartialID.SlateBenefitPartialTransport.endpointFromMassVectorOn_chd

\item The population value of the fixed-support map equals the identified endpoint map. For every array \(b_x\) satisfying
\[
0\le b_x\le m(x)\qquad\text{for all }x,
\]
one has
\[
\sum_{x\in\mathcal X_+(P)}p_xb_x=\sum_xp_xb_x,
\]
because \(x\notin\mathcal X_+(P)\) implies \(p_xm(x)=0\) and hence \(p_xb_x=0\). The capacity validity supplied by \cref{obj:prop:capacity-identification} gives
\[
\ell_i(x)\ge0,\qquad h_j(x)\ge0
\]
for every \(i,j,x\). Hence
\[
q_0(x)=\sum_i\ell_i(x)\ge0,\qquad
q_1(x)=\sum_j h_j(x)\ge0,
\]
so
\[
m(x)=\min\{q_0(x),q_1(x)\}\ge0.
\]
The lower formula in \cref{obj:def:threshold-cuts} includes the candidate \(0\), and therefore
\[
0\le B_L(x).
\]
To prove the upper bound for \(B_L(x)\), it is enough to bound every lower-threshold candidate by \(m(x)\). The zero candidate is bounded by \(m(x)\) by the preceding display. For a threshold \(t\), put
\[
\ell_{\le t}(x)=\sum_{i\le t}\ell_i(x),\qquad
h_{\le t}(x)=\sum_{j\le t}h_j(x),
\]
so that \(0\le h_{\le t}(x)\) and \(\ell_{\le t}(x)\le q_0(x)\). If \(q_0(x)\le q_1(x)\), then \(m(x)=q_0(x)\) and
\[
\ell_{\le t}(x)-h_{\le t}(x)+\min\{\Delta q(x),0\}
\le
\ell_{\le t}(x)
\le q_0(x)=m(x),
\]
because \(-h_{\le t}(x)\le0\) and \(\min\{\Delta q(x),0\}\le0\). If \(q_1(x)\le q_0(x)\), then \(m(x)=q_1(x)\), \(\Delta q(x)=q_1(x)-q_0(x)\le0\), and
\[
\ell_{\le t}(x)-h_{\le t}(x)+\min\{\Delta q(x),0\}
=
\ell_{\le t}(x)-h_{\le t}(x)+q_1(x)-q_0(x)
\le q_1(x)=m(x).
\]
Thus \(B_L(x)\le m(x)\). For the upper threshold cut, \cref{obj:def:threshold-cuts} writes \(B_U(x)\) as the minimum of \(m(x)\) and the threshold sums \(\ell_{<t}(x)+h_{>t}(x)\). The candidate \(m(x)\) gives \(B_U(x)\le m(x)\), while all candidates are nonnegative by capacity validity, so \(0\le B_U(x)\). Consequently
\[
0\le B_L(x)\le m(x),\qquad 0\le B_U(x)\le m(x).
\]
Using the restriction identity with \(b_x=m(x),B_L(x),B_U(x)\) gives
\[
\Phi(\vartheta)
=
\left(
\frac{\sum_xp_xB_L(x)}{M},
\frac{\sum_xp_xB_U(x)}{M}
\right)
=
\Psi(P_{\mathrm{obs}}),
\]
where the final equality is the endpoint formula in \cref{obj:def:identified-interval}. % lean: CausalSmith.PartialID.SlateBenefitPartialTransport.branch_free_pointwise_directional_limit

\item For each cell \(x\), define the empirical screening score
\[
\widehat r_{n,x}
=
\widehat p_{n,x}\widehat m_n(x)
=
\widehat\vartheta_n(X=x)m(\widehat\vartheta_n,x).
\]
The screening rule in \cref{obj:def:plugin-endpoint-estimator} is exactly
\[
x\in\{y:\widehat r_{n,y}>\eta_n\}
\quad\Longleftrightarrow\quad
\eta_n<\widehat r_{n,x}.
\]
If \(p_x>0\), the same positive-denominator argument as above and the directional delta method applied to the scalar score map give tightness of the scaled score error: for every \(\epsilon_{\mathrm{cell}}>0\), there is \(R_x\in\mathbb R\) such that, for all \(n\),
\[
P\!\left(
\max\{R_x,0\}
<
\sqrt n\,|\widehat r_{n,x}-p_xm(x)|
\right)
\le \epsilon_{\mathrm{cell}} .
\]
If \(p_x=0\), then \(P_{\mathrm{obs}}(X=x)=0\), and independent sampling gives
\[
P(\widehat p_{n,x}=0)=1,
\qquad
P(\widehat r_{n,x}=0)=1 .
\]
% lean: CausalSmith.PartialID.SlateBenefitPartialTransport.survivorMassFromMassVector_chd

\item Fix a cell \(x\) and \(\epsilon_{\mathrm{cell}}>0\). If \(p_x=0\), then the preceding display gives exact agreement of the screened and population membership decisions. Suppose \(p_x>0\). If \(x\in\mathcal X_+(P)\), write
\[
r_x=p_xm(x)>0 .
\]
For all sufficiently large \(n\),
\[
\eta_n<r_x/2,
\qquad
\frac{2\max\{R_x,0\}}{r_x}<\sqrt n .
\]
On the event that \(x\) is not screened, \(\widehat r_{n,x}\le\eta_n\), and hence
\[
\sqrt n\,|\widehat r_{n,x}-r_x|
=
\sqrt n\,(r_x-\widehat r_{n,x})
\ge
\sqrt n\,(r_x-\eta_n)
>
\max\{R_x,0\}.
\]
If \(x\notin\mathcal X_+(P)\), then \(p_xm(x)=0\). For all sufficiently large \(n\),
\[
\max\{R_x,0\}<\sqrt n\,\eta_n .
\]
On the event that \(x\) is screened, \(\widehat r_{n,x}>\eta_n\), and therefore
\[
\sqrt n\,|\widehat r_{n,x}-0|
=
\sqrt n\,\widehat r_{n,x}
>
\max\{R_x,0\}.
\]
Thus, for every \(x\),
\[
P\!\left(
\{x\in\{y:\widehat r_{n,y}>\eta_n\}\}
\ne
\{x\in\mathcal X_+(P)\}
\right)
\le \epsilon_{\mathrm{cell}}
\]
eventually. A finite union over \(x\in\mathcal X\) yields
\[
P\!\left(\{x:\widehat r_{n,x}>\eta_n\}=\mathcal X_+(P)\right)\to1.
\]
% lean: CausalSmith.PartialID.SlateBenefitPartialTransport.branch_free_pointwise_directional_limit

\item On the event
\[
\{x:\widehat r_{n,x}>\eta_n\}=\mathcal X_+(P),
\]
the screened sums in \(\widehat\Psi_n\) are exactly the fixed-support sums defining \(\Phi(\widehat\vartheta_n)\). The support is nonempty and each retained score exceeds \(\eta_n>0\), so the screened denominator is positive and the fallback branch in \cref{obj:def:plugin-endpoint-estimator} is inactive on this event. Hence
\[
P\!\left(
\widehat\Psi_n=\Phi(\widehat\vartheta_n)
\right)\to1 .
\]
Together with \(\Phi(\vartheta)=\Psi(P_{\mathrm{obs}})\), equality with probability tending to one transfers the fixed-support weak limit to the screened estimator:
\[
\sqrt n\{\widehat\Psi_n-\Psi(P_{\mathrm{obs}})\}
\rightsquigarrow
D\Psi(\mathbb Z_{P_{\mathrm{obs}}}).
\]
% lean: CausalSmith.PartialID.SlateBenefitPartialTransport.plugInEndpoints_eq_endpointFromMassVectorOn_of_screenedSupport_eq

\item It remains to record consistency. The weak convergence just proved implies tightness of the scaled endpoint error: for every \(\delta>0\) there is \(R_\Psi\in\mathbb R\) such that, for all \(n\),
\[
P\!\left(
R_\Psi<
\sqrt n\,\|\widehat\Psi_n-\Psi(P_{\mathrm{obs}})\|
\right)
\le \delta/2 .
\]
Given \(\epsilon>0\), choose \(n\) large enough that
\[
\frac{\max\{R_\Psi,0\}}{\epsilon}+1\le\sqrt n .
\]
Then
\[
\{\|\widehat\Psi_n-\Psi(P_{\mathrm{obs}})\|>\epsilon\}
\subseteq
\left\{
R_\Psi<
\sqrt n\,\|\widehat\Psi_n-\Psi(P_{\mathrm{obs}})\|
\right\},
\]
and therefore
\[
P\{\|\widehat\Psi_n-\Psi(P_{\mathrm{obs}})\|>\epsilon\}\le\delta/2
\]
eventually. Since \(\delta\) is arbitrary,
\[
\widehat\Psi_n\to_P\Psi(P_{\mathrm{obs}}).
\]
The displayed Gaussian law, covariance identities, support recovery, consistency, and directional weak limit are precisely the asserted conclusions. % lean: CausalSmith.PartialID.SlateBenefitPartialTransport.branch_free_pointwise_directional_limit
\end{enumerate}
\end{proof}

\begin{proof}[Proof of \cref{obj:thm:uniform-deterministic-guard}]
\begin{enumerate}
\item Fix \(n\ge 1\), \(\lambda\in\Lambda\), and \(\omega\).  Let
\(\mathcal Y=\{0,\ldots,K-1\}\).  The finite guard-event collection used here is
\[
\mathcal E
=
\{\mathsf c_x:x\in\mathcal X\}
\cup
\{\mathsf a_{xz}:x\in\mathcal X,\ z\in\{0,1\}\}
\cup
\{\mathsf s_{xz\upsilon i}:x\in\mathcal X,\ z,\upsilon\in\{0,1\},\ i\in\mathcal Y\},
\]
where
\[
\mathsf c_x=\{X=x\},
\qquad
\mathsf a_{xz}=\{X=x,Z=z\},
\qquad
\mathsf s_{xz\upsilon i}=\{X=x,Z=z,D=\upsilon,S=1,\widetilde Y=i\}.
\]
Define
\[
\delta_{n,\lambda}(\omega)
=
\max_{\mathsf e\in\mathcal E}
\left|
\widehat P_{n,\lambda}(\mathsf e;\omega)
-
P_{\mathrm{obs},\lambda}(\mathsf e)
\right|.
\]
The estimator in \cref{obj:def:plugin-endpoint-estimator} uses the totalized
conditional ratio
\[
\operatorname{cond}(u,v)=
\begin{cases}
u/v, & v>0,\\
0, & v\le 0.
\end{cases}
\]
Set
\[
\rho_{\mathrm{cap}}
=
\frac{64K|\mathcal X|}{\varepsilon_Z^2},
\qquad
\mathsf A_n(\delta)
=
\rho_{\mathrm{cap}}\delta+|\mathcal X|\eta_n .
\]
By \cref{obj:def:uniform-law-class}, the law at \(\lambda\) satisfies the
structural restrictions, the i.i.d. sampling condition, and
\cref{obj:ass:uniform-aggregate-survivor-bound}.  Thus
\cref{obj:prop:capacity-identification} gives valid observable capacities and
\(p_x\ge0\).  The overlap condition gives the arm-mass inequality
\[
\varepsilon_Z p_x
\le
P_{\mathrm{obs},\lambda}(X=x,Z=z),
\qquad x\in\mathcal X,
\quad z\in\{0,1\}.
\]
Indeed, when \(p_x=0\) the inequality is immediate from nonnegativity.  When
\(p_x>0\), write
\(\pi(x)=P_{\mathrm{obs},\lambda}(Z=1\mid X=x)\).  For \(z=1\),
\cref{obj:ass:instrument-overlap} gives
\[
P_{\mathrm{obs},\lambda}(X=x,Z=1)
=
 p_x\pi(x)
\ge
\varepsilon_Zp_x .
\]
For \(z=0\), the binary partition and the upper-overlap bound in
\cref{obj:ass:instrument-overlap} give
\[
P_{\mathrm{obs},\lambda}(X=x,Z=0)
=
 p_x\{1-\pi(x)\}
\ge
\varepsilon_Zp_x .
\]
The uniform aggregate bound gives \(m_\star\le M\).

Suppose \(0\le \delta\) and \(\delta_{n,\lambda}(\omega)\le\delta\).  Fix \(x\in\mathcal X\), \(z,\upsilon\in\{0,1\}\), and \(i\in\mathcal Y\).  Write
\[
\mathfrak a_{xz}=P_{\mathrm{obs},\lambda}(X=x,Z=z),
\qquad
\widehat{\mathfrak a}_{xz}
=
\widehat P_{n,\lambda}(X=x,Z=z;\omega),
\]
and
\[
\mathfrak u_{xz\upsilon i}=P_{\mathrm{obs},\lambda}(X=x,Z=z,D=\upsilon,S=1,\widetilde Y=i),
\qquad
\widehat{\mathfrak u}_{xz\upsilon i}
=
\widehat P_{n,\lambda}(X=x,Z=z,D=\upsilon,S=1,\widetilde Y=i;\omega).
\]
The definition of \(\mathcal E\) and the bound
\(\delta_{n,\lambda}(\omega)\le\delta\) give
\[
|\widehat p_x-p_x|\le\delta,
\qquad
|\widehat{\mathfrak a}_{xz}-\mathfrak a_{xz}|\le\delta,
\qquad
|\widehat{\mathfrak u}_{xz\upsilon i}-\mathfrak u_{xz\upsilon i}|\le\delta.
\]
Moreover
\[
0\le \mathfrak u_{xz\upsilon i}\le \mathfrak a_{xz},
\qquad
0\le \widehat{\mathfrak u}_{xz\upsilon i}\le \widehat{\mathfrak a}_{xz},
\]
because each selected-outcome event is contained in its arm event.  Therefore
\[
\left|
\widehat p_x\,
\operatorname{cond}(\widehat{\mathfrak u}_{xz\upsilon i},\widehat{\mathfrak a}_{xz})
-
p_x\,
\operatorname{cond}(\mathfrak u_{xz\upsilon i},\mathfrak a_{xz})
\right|
\le
\frac{8\delta}{\varepsilon_Z^2}.
\]
When \(\delta=0\), the three displayed deviation inequalities force
\(\widehat p_x=p_x\), \(\widehat{\mathfrak a}_{xz}=\mathfrak a_{xz}\), and
\(\widehat{\mathfrak u}_{xz\upsilon i}=\mathfrak u_{xz\upsilon i}\), so the two totalized
conditional expressions are identical, including the case of zero arm mass.
Assume now that \(\delta>0\).  If \(p_x<4\delta/\varepsilon_Z\), both totalized
conditionals lie in \([0,1]\), \(\widehat p_x\le p_x+\delta\), and hence
\[
\left|
\widehat p_x\operatorname{cond}(\widehat{\mathfrak u}_{xz\upsilon i},\widehat{\mathfrak a}_{xz})
-
p_x\operatorname{cond}(\mathfrak u_{xz\upsilon i},\mathfrak a_{xz})
\right|
\le
\widehat p_x+p_x
\le
\frac{8\delta}{\varepsilon_Z^2},
\]
where the last inequality uses \(0<\varepsilon_Z<1/2\).  If
\(p_x\ge4\delta/\varepsilon_Z\), then the arm-overlap inequality gives
\[
4\delta\le \varepsilon_Zp_x\le\mathfrak a_{xz},
\qquad
\widehat{\mathfrak a}_{xz}\ge \mathfrak a_{xz}-\delta\ge\frac34\mathfrak a_{xz}>0.
\]
Thus the ratio branch is active for both denominators, and
\[
\left|
\frac{\widehat{\mathfrak u}_{xz\upsilon i}}{\widehat{\mathfrak a}_{xz}}
-
\frac{\mathfrak u_{xz\upsilon i}}{\mathfrak a_{xz}}
\right|
=
\frac{|\widehat{\mathfrak u}_{xz\upsilon i}\mathfrak a_{xz}-\mathfrak u_{xz\upsilon i}\widehat{\mathfrak a}_{xz}|}
{\widehat{\mathfrak a}_{xz}\mathfrak a_{xz}}
\le
\frac{8\delta}{3\mathfrak a_{xz}}.
\]
The numerator bound follows from
\[
|\widehat{\mathfrak u}_{xz\upsilon i}\mathfrak a_{xz}-\mathfrak u_{xz\upsilon i}\widehat{\mathfrak a}_{xz}|
\le
|\widehat{\mathfrak u}_{xz\upsilon i}-\mathfrak u_{xz\upsilon i}|\mathfrak a_{xz}
+
\mathfrak u_{xz\upsilon i}|\mathfrak a_{xz}-\widehat{\mathfrak a}_{xz}|
\le
2\delta\mathfrak a_{xz},
\]
and the denominator lower bound is
\(\widehat{\mathfrak a}_{xz}\mathfrak a_{xz}\ge3\mathfrak a_{xz}^2/4\).  Since
\(p_x/\mathfrak a_{xz}\le1/\varepsilon_Z\),
\[
\begin{aligned}
&\left|
\widehat p_x\frac{\widehat{\mathfrak u}_{xz\upsilon i}}{\widehat{\mathfrak a}_{xz}}
-
p_x\frac{\mathfrak u_{xz\upsilon i}}{\mathfrak a_{xz}}
\right| \\
&\qquad\le
|\widehat p_x-p_x|\frac{\widehat{\mathfrak u}_{xz\upsilon i}}{\widehat{\mathfrak a}_{xz}}
+
p_x\left|
\frac{\widehat{\mathfrak u}_{xz\upsilon i}}{\widehat{\mathfrak a}_{xz}}
-
\frac{\mathfrak u_{xz\upsilon i}}{\mathfrak a_{xz}}
\right| \\
&\qquad\le
\delta+\frac{8\delta}{3\varepsilon_Z}
\le
\frac{8\delta}{\varepsilon_Z^2},
\end{aligned}
\]
again using \(0<\varepsilon_Z<1/2\).

Each raw lower or upper capacity coordinate is a difference of two such
conditionals, and the projection in \cref{obj:def:plugin-endpoint-estimator} is
coordinatewise positive part, a one-Lipschitz map against the nonnegative
population capacities.  Hence
\[
\max\left\{
\left|\widehat p_x\widehat\ell_i(x)-p_x\ell_i(x)\right|,
\left|\widehat p_x\widehat h_i(x)-p_xh_i(x)\right|
\right\}
\le
\frac{16\delta}{\varepsilon_Z^2}.
\]
Put
\[
e_{\mathrm{cap}}=\frac{16\delta}{\varepsilon_Z^2}.
\]
For each cell define the unconditional population and empirical projected
capacity coordinates by
\[
\mathfrak l_i(x)=p_x\ell_i(x),
\qquad
\mathfrak h_i(x)=p_xh_i(x),
\qquad
\widehat{\mathfrak l}_{i,n}(x)=\widehat p_x\widehat\ell_i(x),
\qquad
\widehat{\mathfrak h}_{i,n}(x)=\widehat p_x\widehat h_i(x).
\]
Let \(\mathfrak m(x)\), \(\mathfrak B_L(x)\), and \(\mathfrak B_U(x)\) denote,
respectively, the survivor-complier mass and the lower and upper threshold cuts
computed from the unconditional capacity array
\((\mathfrak l_i(x),\mathfrak h_i(x))_{i\in\mathcal Y}\).  Let
\(\widehat{\mathfrak m}_n(x)\), \(\widehat{\mathfrak B}_{L,n}(x)\), and
\(\widehat{\mathfrak B}_{U,n}(x)\) denote the corresponding quantities computed
from
\((\widehat{\mathfrak l}_{i,n}(x),\widehat{\mathfrak h}_{i,n}(x))_{i\in\mathcal Y}\).
Since \(p_x\ge0\) and \(\widehat p_x\ge0\),
\cref{obj:lem:weighted-capacity-mass-homogeneity,obj:lem:weighted-capacity-lower-cut-homogeneity,obj:lem:weighted-capacity-upper-cut-homogeneity}
give the homogeneity identities
\[
\mathfrak m(x)=p_xm(x),
\qquad
\mathfrak B_L(x)=p_xB_L(x),
\qquad
\mathfrak B_U(x)=p_xB_U(x),
\]
and
\[
\widehat{\mathfrak m}_n(x)=\widehat p_x\widehat m(x),
\qquad
\widehat{\mathfrak B}_{L,n}(x)=\widehat p_x\widehat B_L(x),
\qquad
\widehat{\mathfrak B}_{U,n}(x)=\widehat p_x\widehat B_U(x).
\]
The preceding coordinate bound says that every lower and upper coordinate in
these unconditional arrays changes by at most \(e_{\mathrm{cap}}\).  Hence each
total \(\sum_i\mathfrak l_i(x)\) or \(\sum_i\mathfrak h_i(x)\) changes by at most
\(Ke_{\mathrm{cap}}\), and the difference of these totals changes by at most
\(2Ke_{\mathrm{cap}}\).  Since \(u\mapsto \min\{u,0\}\), finite maxima, and
finite minima are one-Lipschitz in the sup norm, the lower threshold candidate
\[
\sum_{i\le t}\mathfrak l_i(x)-\sum_{j\le t}\mathfrak h_j(x)
+
\min\left\{\sum_j\mathfrak h_j(x)-\sum_i\mathfrak l_i(x),0\right\}
\]
changes by at most
\[
Ke_{\mathrm{cap}}+Ke_{\mathrm{cap}}+2Ke_{\mathrm{cap}}
=
4Ke_{\mathrm{cap}}.
\]
The zero candidate in the outer maximum is unchanged, so
\(p_xB_L(x)=\mathfrak B_L(x)\) changes by at most \(4Ke_{\mathrm{cap}}\).  For
the upper cut, the mass candidate changes by at most \(Ke_{\mathrm{cap}}\), and
each strict-prefix-plus-tail candidate changes by at most \(2Ke_{\mathrm{cap}}\);
therefore \(p_xB_U(x)=\mathfrak B_U(x)\) changes by at most
\(4Ke_{\mathrm{cap}}\).  The same argument gives the bound
\(4Ke_{\mathrm{cap}}\) for the unconditional survivor mass
\(p_xm(x)=\mathfrak m(x)\).  Summing over the finite covariate support yields
\[
\left|
\sum_x \widehat p_x\widehat m(x)-\sum_x p_xm(x)
\right|
\le
4K|\mathcal X|e_{\mathrm{cap}}
=
\rho_{\mathrm{cap}}\delta,
\]
and
\[
\left|
\sum_x \widehat p_x\widehat B_L(x)-\sum_x p_xB_L(x)
\right|
\le
\rho_{\mathrm{cap}}\delta,
\qquad
\left|
\sum_x \widehat p_x\widehat B_U(x)-\sum_x p_xB_U(x)
\right|
\le
\rho_{\mathrm{cap}}\delta.
\]
Before applying the screening rule, record the cut bounds used for both the
population capacities and the projected empirical capacities.  Fix a covariate
cell \(x\) and either one of these two nonnegative capacity arrays; write its
associated totals, gap, mass, and cuts with a superscript \(\circ\):
\[
q_0^\circ(x)=\sum_{i\in\mathcal Y}\ell_i^\circ(x),
\qquad
q_1^\circ(x)=\sum_{j\in\mathcal Y}h_j^\circ(x),
\]
\[
\Delta q^\circ(x)=q_1^\circ(x)-q_0^\circ(x),
\qquad
m^\circ(x)=\min\{q_0^\circ(x),q_1^\circ(x)\}.
\]
All coordinates are nonnegative, so \(m^\circ(x)\ge0\).  By the formula in
\cref{obj:def:threshold-cuts}, \(B_L^\circ(x)\) is the maximum of the zero
candidate and the threshold candidates, and therefore
\[
0\le B_L^\circ(x).
\]
The upper cut is the minimum of \(m^\circ(x)\) and nonnegative
strict-prefix-plus-tail sums, so
\[
0\le B_U^\circ(x)
\qquad\text{and}\qquad
B_U^\circ(x)\le m^\circ(x).
\]
It remains to bound the lower cut above by the mass.  For a threshold
\(t\in\mathcal Y\), define its lower-cut candidate by
\[
C_t^\circ(x)
=
\ell_{\le t}^\circ(x)-h_{\le t}^\circ(x)
+\min\{\Delta q^\circ(x),0\}.
\]
If \(q_0^\circ(x)\le q_1^\circ(x)\), then
\(m^\circ(x)=q_0^\circ(x)\), \(h_{\le t}^\circ(x)\ge0\), and
\(\min\{\Delta q^\circ(x),0\}\le0\); hence
\[
C_t^\circ(x)
\le
\ell_{\le t}^\circ(x)
\le
q_0^\circ(x)
=
m^\circ(x).
\]
If \(q_1^\circ(x)\le q_0^\circ(x)\), then
\(m^\circ(x)=q_1^\circ(x)\) and
\(\min\{\Delta q^\circ(x),0\}=q_1^\circ(x)-q_0^\circ(x)\), so
\[
C_t^\circ(x)
=
\ell_{\le t}^\circ(x)-h_{\le t}^\circ(x)
+q_1^\circ(x)-q_0^\circ(x)
\le
q_1^\circ(x)
=
m^\circ(x),
\]
because \(\ell_{\le t}^\circ(x)\le q_0^\circ(x)\) and
\(h_{\le t}^\circ(x)\ge0\).  The zero candidate is also at most
\(m^\circ(x)\).  Thus every candidate in the finite maximum defining the lower
cut is at most \(m^\circ(x)\), and
\[
B_L^\circ(x)\le m^\circ(x).
\]
Applying these inequalities first to the valid population capacities and then
to the projected empirical capacities gives, for every \(x\),
\[
0\le B_L(x),B_U(x)\le m(x),
\qquad
0\le \widehat B_L(x),\widehat B_U(x)\le \widehat m(x).
\]
Screening removes, in each cell, at most \(\eta_n\) from the empirical mass and
at most \(\eta_n\) from each empirical numerator.  Indeed, a screened-out cell
satisfies \(\widehat p_x\widehat m(x)\le\eta_n\), while the just-proved
empirical cut bounds give
\(0\le \widehat B_L(x)\le\widehat m(x)\) and
\(0\le\widehat B_U(x)\le\widehat m(x)\).  Thus, with
\[
N_L=\sum_x p_xB_L(x),
\qquad
N_U=\sum_x p_xB_U(x),
\]
\[
\left|\widehat M_n-M\right|\le \mathsf A_n(\delta),
\qquad
\left|\widehat N_{L,n}-N_L\right|\le \mathsf A_n(\delta),
\qquad
\left|\widehat N_{U,n}-N_U\right|\le \mathsf A_n(\delta).
\]
Since \(p_x\ge0\), \(\widehat p_x\ge0\), and the retained-cell indicators are
nonnegative, the same cut bounds also imply
\[
0\le N_L,N_U\le M,
\qquad
0\le\widehat N_{L,n},\widehat N_{U,n}\le\widehat M_n.
\]

If \(4\mathsf A_n(\delta)<m_\star\), then
\[
\widehat M_n
\ge
M-\mathsf A_n(\delta)
\ge
m_\star-\mathsf A_n(\delta)
>
\frac{3m_\star}{4},
\]
so the ratio branch of \(\widehat\Psi_n\) is active.  For either endpoint
numerator \(N\in\{N_L,N_U\}\) and its empirical screened counterpart
\(\widehat N\),
\[
\left|
\frac{\widehat N}{\widehat M_n}-\frac{N}{M}
\right|
\le
\frac{|\widehat N-N|}{\widehat M_n}
+
\frac{N|M-\widehat M_n|}{M\widehat M_n}
\le
\frac{2\mathsf A_n(\delta)}{\widehat M_n}
\le
\frac{3\mathsf A_n(\delta)}{m_\star}.
\]
If \(4\mathsf A_n(\delta)\ge m_\star\), then
\(\min\{1,4\mathsf A_n(\delta)/m_\star\}=1\), and both the target endpoints and
the plug-in endpoints lie in \([0,1]\).  Combining the two cases gives
\[
\max\left\{
\left|\widehat\theta_{L,n,\lambda}(\omega)-\theta_{L,\lambda}\right|,
\left|\widehat\theta_{U,n,\lambda}(\omega)-\theta_{U,\lambda}\right|
\right\}
\le
\min\left\{1,\frac{4\mathsf A_n(\delta)}{m_\star}\right\},
\]
and the sharper bound
\[
4\mathsf A_n(\delta)<m_\star
\quad\Longrightarrow\quad
\max\left\{
\left|\widehat\theta_{L,n,\lambda}(\omega)-\theta_{L,\lambda}\right|,
\left|\widehat\theta_{U,n,\lambda}(\omega)-\theta_{U,\lambda}\right|
\right\}
\le
\frac{3\mathsf A_n(\delta)}{m_\star}.
\]
% lean: CausalSmith.PartialID.SlateBenefitPartialTransport.plugInEndpoints_error_le_of_maxDeviation

\item For \(\tau>0\), the maximal guard deviation has the finite second-moment
tail bound
\[
\mu_\lambda\{\omega:\tau<\delta_{n,\lambda}(\omega)\}
\le
\frac{|\mathcal E|}{4n\tau^2}.
\]
To see this, for \(\mathsf e\in\mathcal E\) put
\[
W_{r,\mathsf e}
=
\mathbf 1\{\mathsf e(O_{\lambda,r})\}
-
P_{\mathrm{obs},\lambda}(\mathsf e),
\qquad r=1,\ldots,n .
\]
The i.i.d. sampling condition gives independent mean-zero variables with
\[
\operatorname{Var}_{\mu_\lambda}(W_{r,\mathsf e})
=
P_{\mathrm{obs},\lambda}(\mathsf e)
\{1-P_{\mathrm{obs},\lambda}(\mathsf e)\}
\le \frac14.
\]
Hence
\[
\widehat P_{n,\lambda}(\mathsf e;\omega)
-
P_{\mathrm{obs},\lambda}(\mathsf e)
=
\frac1n\sum_{r=1}^n W_{r,\mathsf e}(\omega),
\qquad
\operatorname{Var}_{\mu_\lambda}\!\left(
\frac1n\sum_{r=1}^n W_{r,\mathsf e}
\right)
\le
\frac{1}{4n}.
\]
Chebyshev's inequality and a finite union bound over \(\mathcal E\) give the
displayed tail bound.
% lean: CausalSmith.PartialID.SlateBenefitPartialTransport.maxDeviation_tail_le
Since \(\mathcal X\) is nonempty, \(\mathcal E\) contains
a cell event, so
\[
b_{n,\alpha}
=
\sqrt{\frac{|\mathcal E|}{4n\alpha}}>0
\]
and therefore
\[
\mu_\lambda\{\omega:b_{n,\alpha}<\delta_{n,\lambda}(\omega)\}
\le
\alpha .
\]
% lean: CausalSmith.PartialID.SlateBenefitPartialTransport.maxDeviation_unionThreshold_tail_le

\item On the event
\(\{\omega:\delta_{n,\lambda}(\omega)\le b_{n,\alpha}\}\), the first step with
\(\delta=b_{n,\alpha}\) and \cref{obj:def:face-aware-inference-handle} give
\[
\max\left\{
\left|\widehat\theta_{L,n,\lambda}(\omega)-\theta_{L,\lambda}\right|,
\left|\widehat\theta_{U,n,\lambda}(\omega)-\theta_{U,\lambda}\right|
\right\}
\le
g_{n,\alpha}.
\]
Let \(y\in\Theta_I(P_{\mathrm{obs},\lambda})\).  By
\cref{obj:def:identified-interval},
\[
\theta_{L,\lambda}\le y\le \theta_{U,\lambda},
\]
and the endpoint bounds established from \(0\le B_L(x),B_U(x)\le m(x)\) give
\(0\le\theta_{L,\lambda}\le1\) and \(0\le\theta_{U,\lambda}\le1\).  Thus
\[
\max\{0,\widehat\theta_{L,n,\lambda}(\omega)-g_{n,\alpha}\}
\le y
\le
\min\{1,\widehat\theta_{U,n,\lambda}(\omega)+g_{n,\alpha}\},
\]
which is exactly
\(y\in\mathcal C_{1-\alpha,n,\lambda}(\omega)\).  The tail bound in the previous
step therefore implies
\[
\mu_\lambda\!\left\{
\omega:
\Theta_I(P_{\mathrm{obs},\lambda})
\subseteq
\mathcal C_{1-\alpha,n,\lambda}(\omega)
\right\}
\ge
1-\alpha .
\]
Taking the infimum over \(\lambda\in\Lambda\) proves the finite-sample coverage
claim for this \(n\), and \(n\ge1\) was arbitrary.
% lean: CausalSmith.PartialID.SlateBenefitPartialTransport.uniform_deterministic_guard

\item Define
\[
\tau_n
=
\frac{|\mathcal X|}{4\rho_{\mathrm{cap}}}\eta_n .
\]
Since \(K\ge3\), \(|\mathcal X|>0\), and \(\varepsilon_Z>0\), we have
\(\rho_{\mathrm{cap}}>0\).  Together with the screening-threshold assumption
\(\eta_n>0\), this gives \(\tau_n>0\).  Moreover,
\[
4\{\rho_{\mathrm{cap}}\tau_n+|\mathcal X|\eta_n\}
=
5|\mathcal X|\eta_n
\longrightarrow 0,
\]
so eventually
\[
4\{\rho_{\mathrm{cap}}\tau_n+|\mathcal X|\eta_n\}<m_\star .
\]
For such \(n\), if
\(\delta_{n,\lambda}(\omega)\le\tau_n\), the sharper deterministic bound from
the first step gives
\[
\max\left\{
\left|\widehat\theta_{L,n,\lambda}(\omega)-\theta_{L,\lambda}\right|,
\left|\widehat\theta_{U,n,\lambda}(\omega)-\theta_{U,\lambda}\right|
\right\}
\le
\frac{3\{\rho_{\mathrm{cap}}\tau_n+|\mathcal X|\eta_n\}}{m_\star}.
\]
Because \(\rho_{\mathrm{cap}}\tau_n=|\mathcal X|\eta_n/4\) and
\(b_{n,\alpha}\ge0\),
\[
\frac{3\{\rho_{\mathrm{cap}}\tau_n+|\mathcal X|\eta_n\}}{m_\star}
\le
\frac{4\{\rho_{\mathrm{cap}}b_{n,\alpha}+|\mathcal X|\eta_n\}}{m_\star}.
\]
Together with the bound by \(1\) from the first deterministic inequality, this
yields
\[
\max\left\{
\left|\widehat\theta_{L,n,\lambda}(\omega)-\theta_{L,\lambda}\right|,
\left|\widehat\theta_{U,n,\lambda}(\omega)-\theta_{U,\lambda}\right|
\right\}
\le
g_{n,\alpha}.
\]
Consequently, eventually in \(n\),
\[
\left\{\omega:
g_{n,\alpha}
<
\max\left(
\left|\widehat\theta_{L,n,\lambda}(\omega)-\theta_{L,\lambda}\right|,
\left|\widehat\theta_{U,n,\lambda}(\omega)-\theta_{U,\lambda}\right|
\right)
\right\}
\subseteq
\{\omega:\tau_n<\delta_{n,\lambda}(\omega)\}.
\]
The tail bound gives, uniformly in \(\lambda\),
\[
\mu_\lambda\{\omega:\tau_n<\delta_{n,\lambda}(\omega)\}
\le
\frac{|\mathcal E|}{4n\tau_n^2}
=
\frac{4\rho_{\mathrm{cap}}^2|\mathcal E|}
{n|\mathcal X|^2\eta_n^2}.
\]
% lean: CausalSmith.PartialID.SlateBenefitPartialTransport.maxDeviation_tail_le
The right-hand side tends to \(0\) because
\(\sqrt n\,\eta_n\to\infty\).
% lean: CausalSmith.PartialID.SlateBenefitPartialTransport.inverse_sqrtEta_tail_tendsto_zero
This proves
\[
\sup_{\lambda\in\Lambda}
\mu_\lambda\!\left\{
\omega:
g_{n,\alpha}
<
\max\left\{
\left|\widehat\theta_{L,n,\lambda}(\omega)-\theta_{L,\lambda}\right|,
\left|\widehat\theta_{U,n,\lambda}(\omega)-\theta_{U,\lambda}\right|
\right\}
\right\}
\longrightarrow 0 .
\]
% lean: CausalSmith.PartialID.SlateBenefitPartialTransport.uniform_deterministic_guard

\item Let
\[
\beta_\alpha
=
\sqrt{\frac{|\mathcal E|}{4\alpha}} .
\]
For \(n\ge1\),
\[
b_{n,\alpha}=\beta_\alpha n^{-1/2}.
\]
Therefore \cref{obj:def:face-aware-inference-handle} gives
\[
0\le g_{n,\alpha}
\le
\frac{4}{m_\star}
\left\{
\rho_{\mathrm{cap}}\beta_\alpha n^{-1/2}
+
|\mathcal X|\eta_n
\right\}
\le
\frac{4(\rho_{\mathrm{cap}}\beta_\alpha+|\mathcal X|)}{m_\star}
\left(n^{-1/2}+\eta_n\right),
\]
which proves
\[
g_{n,\alpha}=O(n^{-1/2}+\eta_n).
\]

Now let \(L_n>0\), \(L_n\to\infty\), and \(L_n=o(\sqrt n)\), and set
\[
\eta_n=n^{-1/2}L_n .
\]
For every \(n\ge1\), \(\eta_n>0\).  Also
\[
\eta_n=\frac{L_n}{\sqrt n}\longrightarrow0,
\qquad
\sqrt n\,\eta_n=L_n\longrightarrow\infty .
\]
Since \(L_n\to\infty\), eventually \(L_n\ge1\), and hence
\(n^{-1/2}\le n^{-1/2}L_n=\eta_n\).  The preceding guard bound then gives,
eventually,
\[
g_{n,\alpha}
\le
\frac{8(\rho_{\mathrm{cap}}\beta_\alpha+|\mathcal X|)}{m_\star}
\,n^{-1/2}L_n,
\]
so
\[
g_{n,\alpha}=O(n^{-1/2}L_n).
\]
% lean: CausalSmith.PartialID.SlateBenefitPartialTransport.uniform_deterministic_guard
\end{enumerate}
\end{proof}
